\documentclass[pdflatex,sn-mathphys-num]{sn-jnl}

\usepackage{graphicx}% Include figure files
\usepackage{dcolumn}% Align table columns on decimal point
\usepackage{bm}% bold math
\usepackage{enumitem}
\usepackage{graphicx}%
\usepackage{multirow}%
\usepackage{amsmath,amssymb,amsfonts}%
\usepackage{amsthm}%
\usepackage{mathrsfs}%
\usepackage[title]{appendix}%
\usepackage{xcolor}%
\usepackage{textcomp}%
\usepackage{booktabs}%
\usepackage{algorithm}%
\usepackage{algorithmicx}%
\usepackage{algpseudocode}%
\usepackage{listings}
\usepackage{aliascnt}
\usepackage{slashed}
\usepackage{hyperref}% add hypertext capabilities
%\usepackage[mathlines]{lineno}% Enable numbering of text and display math
%\linenumbers\relax % Commence numbering lines


%% as per the requirement new theorem styles can be included as shown below
\theoremstyle{thmstyleone}
\newtheorem{theorem}{Theorem}

% Lemma
\newaliascnt{lemma}{theorem}
\newtheorem{lemma}[lemma]{Lemma}
\aliascntresetthe{lemma}
\providecommand*{\lemmaautorefname}{Lemma}

% Axiom
\newaliascnt{axiom}{theorem}
\newtheorem{axiom}[axiom]{Axiom}
\aliascntresetthe{axiom}
\providecommand*{\axiomautorefname}{Axiom}

% Corollary
\newaliascnt{corollary}{theorem}
\newtheorem{corollary}[corollary]{Corollary}
\aliascntresetthe{corollary}
\providecommand*{\corollaryautorefname}{Corollary}

% Proposition
\newaliascnt{proposition}{theorem}
\newtheorem{proposition}[proposition]{Proposition}
\aliascntresetthe{proposition}
\providecommand*{\propositionautorefname}{Proposition}

\theoremstyle{thmstylethree}%

% Definition
\newaliascnt{definition}{theorem}
\newtheorem{definition}[definition]{Definition}
\aliascntresetthe{definition}
\providecommand*{\definitionautorefname}{Definition}


\theoremstyle{thmstyletwo}%

\newaliascnt{example}{theorem}
\newtheorem{example}[example]{Example}
\aliascntresetthe{example}
\providecommand*{\exampleautorefname}{Example}

\newaliascnt{remark}{theorem}
\newtheorem{remark}[remark]{Remark}
\aliascntresetthe{remark}
\providecommand*{\remarkautorefname}{Remark}

%\usepackage[showframe,%Uncomment any one of the following lines to test 
%%scale=0.7, marginratio={1:1, 2:3}, ignoreall,% default settings
%%text={7in,10in},centering,
%%margin=1.5in,
%%total={6.5in,8.75in}, top=1.2in, left=0.9in, includefoot,
%%height=10in,a5paper,hmargin={3cm,0.8in},
%]{geometry}


\begin{document}

\newcommand{\SU}[1]{\mathrm{SU(#1)}}
\newcommand{\SMSubgroup}{\SU{3}\times\SU{2}\times\U{1}}
\newcommand{\SO}[1]{\mathrm{SO(#1)}}
\newcommand{\U}[1]{\mathrm{U(#1)}}
\newcommand{\E}[1]{\mathrm{E_{#1}}}
\newcommand{\M}{\mathcal{M}}
\newcommand{\Oct}{\mathbb{O}}
\newcommand{\jalg}{J_3(\mathbb{O})}
\newcommand{\SOTen}{\SO{10}}
\newcommand{\UOne}{\U{1}}
\newcommand{\ESix}{\E{6}}
\newcommand{\SOTenXUOne}{\SOTen \times \UOne}
\newcommand{\MCoset}{\frac{\ESix}{\SOTenXUOne}}
\newcommand{\MCosetInline}{\ESix/(\SOTenXUOne)}
\newcommand{\ESixBreaking}{\ESix \;\to\; \SOTenXUOne}
\newcommand{\map}[3]{#1 \ : \ \#2 \;\to\; \#3}
\newcommand{\set}[2]{\{ #1 \,:\, #2 \}}
\newcommand{\Herm}[2]{\mathrm{Herm}_{#1}(#2)}
\newcommand{\chr}[1]{\mathrm{char}(#1)}
\newcommand{\R}{\mathbb{R}}
\newcommand{\C}{\mathbb{C}}
\newcommand{\Tr}{\operatorname{Tr}}
\newcommand{\s}{\hspace{.1em}}
\newcommand{\sm}{\hspace{.065em}}
\newcommand{\Str}{\operatorname{Str}}
\newcommand{\Spin}  {\operatorname{Spin}}
\newcommand{\smexp}[1]{^{\sm #1}}
\newcommand{\Hess}[1]{\operatorname{Hess}}
\newcommand{\dv}[1]{\mathrm{d}\smexp{#1}}

\title[Universal coset field theory]{Universal coset field theory}

\author*[1]{\fnm{Brandon} \sur{Belna}}\email{bbelna@iu.edu}
\affil*[1]{\orgdiv{Department of Mathematics}, \orgname{Indiana University East}, \orgaddress{\street{2325 Chester Blvd}, \city{Richmond}, \postcode{47374}, \state{IN}, \country{USA}}}

\abstract{
We formulate a universal coset field theory (UCFT), a four-dimensional
Euclidean quantum field theory whose exceptional datum is the Albert algebra
$\jalg$ and the associated compact Lie type $\ESix$. Its central result is
that, under the two stated axioms, gravity is not a fundamental interaction:
spacetime and the metric that propagates on it are emergent, induced
structures rather than postulated inputs, and no fundamental graviton field
ever appears. An $\ESix$-invariant coercive
potential has vacuum orbit Cartan's Hermitian symmetric space
EIII, $\MCosetInline$; its $32$ real tangent directions form the irreducible
$\SOTenXUOne$ multiplet with complexification
$\mathbf{16}_{-3}\oplus\overline{\mathbf{16}}_{+3}$ and carry a unique
positive invariant metric. For nondegenerate spacetime-filling maps into EIII,
the pull-back of this metric gives a Riemannian first fundamental form, so the
spacetime metric is derived from the coset field configuration rather than
assumed \emph{a priori}. Lorentzian physics is interpreted through
Osterwalder--Schrader reconstruction, and one-loop matter determinants induce
an Einstein--Hilbert term in the sense of Sakharov: gravity is thus a
radiative effect of matter loops on the coset, not a fundamental field, at
every stage of the construction. We compare this coset sector with the gauge-unfixed Higgs orbit of an
$E_{8}\times E_{8}$ heterotic compactification with visible $\ESix$. At the level
of the two-derivative bosonic action and the $\SOTenXUOne_\psi$ current
algebra, the visible sector matches; a $\mathbb Z_{12}$ asymmetric
orbifold with three generations and an adjoint Higgs supplies the candidate
realization. The resulting Albert--EIII criterion is relative: within the
Albert-selected heterotic class, a theory can realize the correspondence
only if its residual Higgs/moduli sector is the EIII coset with the stated
tangent representation, metric, gauge data, anomaly and Green--Schwarz
matching, and worldsheet consistency. In the stated truncation, the functional
renormalization-group flow has a unique ultraviolet-attractive non-Gaussian
fixed point and gives
$M_{\mathrm{GUT}}\simeq2\times10\smexp{16}\,\mathrm{GeV}$. We also show that the
Riemann hypothesis is equivalent to the absence of complex mass-squared modes,
which is a spectral reformulation of the prime-orbit trace of the finite adelic
zeta sector.
We distinguish the representation-theoretic consequences from assumptions
requiring further dynamical input.
}

\keywords{swampland, exceptional Jordan algebra, $E_6$, heterotic string,
$\mathbb Z_{12}$ orbifold, coset field theory, EIII, induced gravity, emergent
spacetime, asymptotic safety, grand unification, Riemann zeta function,
spectral determinant}

%%\pacs[JEL Classification]{D8, H51}

%%\pacs[MSC Classification]{35A01, 65L10, 65L12, 65L20, 65L70}

\maketitle

%======================================================================%
\section{Introduction}
%======================================================================%

Exceptional algebraic structures recur throughout high-energy physics. The
exceptional Jordan (Albert) algebra $\jalg$ of $3\times3$ Hermitian octonionic
matrices has automorphism group $F_{4}$ and reduced structure group of complex
Lie type $E_{6}$, the type that preserves its cubic norm
\cite{Jacobson:1968,Springer:2000,McCrimmon:2004}. In this paper
$\ESix$ denotes the compact real form used in the EIII coset construction. It
contains $\SOTenXUOne$, under which the fundamental $\mathbf{27}$ branches as
\[
  \mathbf{27}=\mathbf{16}_{+1}\oplus\mathbf{10}_{-2}\oplus\mathbf{1}_{+4},
\]
the $\mathbf{16}$ of $\Spin(10)$ carrying one Standard-Model family together
with a right-handed neutrino \cite{Slansky:1981}. Formally real Jordan algebras
were introduced as algebras of quantum observables; $\jalg$ is the unique
exceptional member of their classification. The organizing question is how much
four-dimensional field theory is fixed once this exceptional datum and its
compact $E_{6}$ symmetry are specified.

The main result of this paper answers a stronger question than that
organizing one: once the Albert datum and its two axioms are fixed, gravity
is not among the fundamental ingredients of the theory. No graviton field, no
spacetime metric, and no gravitational action are postulated anywhere in the
construction. Spacetime itself is the image of a spacetime-filling map into
the EIII coset, its metric is the pull-back of the unique coset metric
(\autoref{thm:EmergentRiemannian}), and the Einstein--Hilbert term that governs
its dynamics is generated radiatively, at one loop, by the matter content the
same construction forces into existence (\autoref{thm:inducedEH}). In this
precise sense, \emph{if UCFT holds, gravity is not fundamental}: it is a
low-energy, induced consequence of the Albert-algebraic symmetry breaking, on
the same footing as the Sakharov--Adler mechanism, rather than an independent
input to the theory.

The construction rests on two axioms. The first posits a maximally symmetric initial state
in $\jalg$; the second posits the usual quantum-mechanical kinematics of a
universal wavefunction with unitary evolution. The field theory is then defined
by an $\ESix$-invariant symmetry-breaking sector on the Jordan state space. A
polynomial potential, built from the quadratic and cubic invariants $Q$ and
$N$, is coercive and is minimized on the selected orbit
\[
  \ESix\cdot X_{0}\cong\MCoset,
\]
the compact Hermitian symmetric space EIII \cite{Helgason:1978}, of real
dimension $78-46=32$. The $32$ Goldstone directions form a single irreducible
multiplet of $H=\SOTenXUOne$, whose complexification is
$\mathbf{16}_{-3}\oplus\overline{\mathbf{16}}_{+3}$; the invariant metric is
therefore unique up to scale and positive definite.

The coset sector also supplies the geometric data. A spacetime-filling
configuration of the Goldstone fields immerses a four-manifold into EIII; its
first fundamental form is a nondegenerate Riemannian metric, so the
four-dimensional geometry is induced by the field configuration rather than
postulated independently. Since the functional integral is Euclidean, physical
Lorentzian spacetime is understood through Osterwalder--Schrader reconstruction
or, equivalently, analytic continuation. Quantizing the coset scalars, gauge
fields, and a chiral $\mathbf{16}$ induces an Einstein--Hilbert term at one
loop in the sense of Sakharov
\cite{SakharovInducedGravity,AdlerInducedGravity}; the metric enters through
this induced action rather than through a fundamental graviton.

Two additional ingredients are analyzed. First, within the stated functional
renormalization-group truncation, the coset-field action has an
ultraviolet-attractive non-Gaussian fixed point and a grand-unified matching
scale $M_{\mathrm{GUT}}\simeq2\times10\smexp{16}\,\mathrm{GeV}$. Second, the
$\ESix$ branchings supply a grand-unified matter skeleton and anomaly
constraints. We also include an arithmetic extension: a finite adelic zeta
sector whose prime-orbit trace reproduces the logarithmic derivative of the
completed Riemann zeta factor, and whose determinant realization gives a
mass-spectrum reformulation of the Riemann hypothesis.

The paper is organized as follows. \autoref{sec:AxiomsAlbertSSB} states the
axioms, recalls the Albert algebra and its $\ESix$ symmetry, and constructs the
symmetry-breaking potential and its vacuum manifold.
\autoref{sec:EmergentFields} derives the emergent four-dimensional geometry and
the low-energy field content. \autoref{sec:sigma} defines the coset
$\sigma$-model, quantizes it, and derives the induced Einstein--Hilbert term and
the coset mass spectrum. \autoref{sec:FRG} presents the renormalization-group
flow, the non-Gaussian fixed point, and the grand-unified matching scale.
\autoref{sec:ZetaSector} adds the arithmetic zeta sector and the
spectral reformulation of the Riemann hypothesis. \autoref{sec:Holographic}
compares the coset sector with a heterotic visible sector, \autoref{sec:Z12}
records the $\mathbb Z_{12}$ candidate realization, and
\autoref{sec:Swampland} formulates the resulting relative swampland criterion.
Proofs and technical computations are collected in the appendices.

%======================================================================%
\section{Axioms, the Albert Algebra, and Spontaneous Symmetry Breaking}
\label{sec:AxiomsAlbertSSB}
%======================================================================%
This section states the axioms, recalls the Albert algebra and the compact
$E_{6}$ symmetry used in the coset construction, constructs the algebraic
symmetry-breaking potential, and counts the resulting $32$ Goldstone
directions.

%%%%%%%%%%%%%%%%%%%%%%%%%%%%%%%%%%%%%%%%%%%%%%%%%%%%%%%%%%%%%%%%%%%%%%
\subsection{Axioms}\label{subsec:Axioms}
%%%%%%%%%%%%%%%%%%%%%%%%%%%%%%%%%%%%%%%%%%%%%%%%%%%%%%%%%%%%%%%%%%%%%%
\begin{axiom}[Maximally symmetric vacuum]\label{ax:AxiomI}
There is a distinguished state $\Phi\in\jalg$ fixed by the full exceptional
symmetry selected by the Albert algebra.
\end{axiom}
\begin{axiom}[Universal wavefunction]\label{ax:AxiomII}
Physical states $\psi\in\mathcal H$ evolve through a one-parameter family
$\{U(t)\}_{t\in\mathbb R}$ of unitary operators satisfying
$U(0)=\mathrm{Id}$, $U(t+s)=U(t)U(s)$, and strong continuity.
\end{axiom}

\begin{remark}
\autoref{ax:AxiomII} provides the kinematical Hilbert space and the
one-parameter family of unitary operators. Whenever we later write a functional
integral, we treat it as a formal generator of connected Green functions; its
constructive interpretation is through the Osterwalder--Schrader
reconstruction theorem applied to Euclidean correlation functions
\cite{GlimmJaffe:1987}.
\end{remark}


No additional axioms are introduced; all further ingredients are definitions,
constructions, or consequences within the framework specified by these two
axioms.

%%%%%%%%%%%%%%%%%%%%%%%%%%%%%%%%%%%%%%%%%%%%%%%%%%%%%%%%%%%%%%%%%%%%%%
\subsection{The Albert Algebra}
\label{subsec:Albert}
%%%%%%%%%%%%%%%%%%%%%%%%%%%%%%%%%%%%%%%%%%%%%%%%%%%%%%%%%%%%%%%%%%%%%%

\begin{definition}[Jordan algebra {\cite{Jacobson:1968}}]
Let $F$ be a field with $\operatorname{char} F \neq 2$.  
A \emph{Jordan algebra} is a commutative algebra $(A,\circ)$ satisfying
\[
  a\circ(b\circ a\smexp{2}) \;=\;(a\circ b)\circ a\smexp{2},
  \quad a,b\in A.
\]
\end{definition}

\begin{definition}[Albert algebra {\cite{Albert:1934,McCrimmon:2004}}]
The \emph{Albert algebra} is
\[
  \jalg=\{X\in M_{3}(\Oct) : X=X\smexp{\dagger}\},
  \quad
  X\circ Y=\frac12\,(XY+YX).
\]
\end{definition}

\begin{theorem}\label{thm:JordanDim}
The Albert algebra is a $27$-dimensional real Jordan algebra.
\end{theorem}

\begin{theorem}[Quadratic and cubic invariants]\label{thm:Invs}
$\jalg$ admits  
\[
  Q(X)=\Tr(X\smexp{2}), 
  \quad
  N(X)=\det_{\Oct}(X),
\]
which are preserved by the corresponding reduced structure group.
\end{theorem}

\begin{theorem}[Exceptional type {\cite{Springer:2000}}]\label{thm:E6}
The reduced structure algebra of $\jalg$ has complexification of type
$\mathfrak e_{6}$. In the compact coset construction below, $\ESix$ denotes the
simply connected compact real form of this Lie type.
\end{theorem}

\begin{theorem}[Regular orbit criterion]\label{thm:OrbitClass}
On the regular compact stratum used below, the invariants $Q$ and $N$ separate
$\ESix$-orbits: for $X,Y$ in that stratum,
\[
  \bigl(Q(X),N(X)\bigr)=\bigl(Q(Y),N(Y)\bigr)
\]
if and only if there exists $g\in\ESix$ such that $g\cdot X = Y$.
\end{theorem}

\begin{corollary}[Existence of an $\SOTenXUOne$ stabilizer]
\label{cor:Isotropy}
There exists $X_{0}\in J_{3}(\mathbb{O})$ with
$\mathrm{Stab}_{\ESix}(X_{0})\cong \SOTen\times \UOne$.
\end{corollary}

%%%%%%%%%%%%%%%%%%%%%%%%%%%%%%%%%%%%%%%%%%%%%%%%%%%%%%%%%%%%%%%%%%%%%%
\subsection{Polynomial Symmetry-Breaking Potential}\label{subsec:PolyPot}
%%%%%%%%%%%%%%%%%%%%%%%%%%%%%%%%%%%%%%%%%%%%%%%%%%%%%%%%%%%%%%%%%%%%%%
Set
\[
  I_{2}(X)=Q(X),
  \quad
  I_{3}(X)=N(X),
  \quad
  c_{2}=I_{2}(X_{0}),
  \quad
  c_{3}=I_{3}(X_{0}).
\]
We choose $X_{0}$ in the regular stratum of \autoref{thm:OrbitClass}. Then the
simultaneous level set through $X_{0}$ is the orbit
$\ESix\cdot X_{0}$. Define
\begin{equation}\label{eq:Pdef}
\begin{aligned}
  P(X)
  &=\bigl[I_{2}(X)-c_{2}\bigr]\smexp{2}
  +\bigl[I_{3}(X)-c_{3}\bigr]\smexp{2},\\[2pt]
  V(X)
  &=P(X)
  +\alpha\,\bigl[I_{2}(X)-c_{2}\bigr]\smexp{k},
\end{aligned}
\end{equation}
with $\alpha>0$ and even $k\ge2$.

\begin{lemma}\label{lem:Zeros}
$P(X)=0$ if and only if $X\in\ESix\cdot X_{0}$.
\end{lemma}

\begin{proof}
By invariance of $I_{2}$ and $I_{3}$ under $\ESix$,
$P(g\cdot X_{0})=0$. Conversely, if $P(X)=0$, then
$(I_{2},I_{3})=(c_{2},c_{3})$, so \autoref{thm:OrbitClass} gives
$X\in\ESix \cdot X_{0}$.
\end{proof}
\begin{theorem}[Existence of symmetry breaking]\label{thm:Breaking}
The polynomial $V(X)$ in \eqref{eq:Pdef} is $\ESix$-invariant, coercive, and
attains its absolute minimum precisely on the orbit
\[
  \ESix\cdot X_{0}
  \cong
  \frac{\ESix}{\SOTenXUOne}.
\]
\end{theorem}

\begin{proof}
Invariance is immediate from the invariance of $I_{2}$ and $I_{3}$. Since
$I_{2}(X)$ is positive definite on $\jalg$, the even-degree term
$\alpha[I_{2}(X)-c_{2}]\smexp{k}$ dominates as $\|X\|\to\infty$, so $V$ is
coercive. By \autoref{lem:Zeros}, $P\ge0$ with $P=0$ exactly on
$\ESix\cdot X_{0}$; the additional term is nonnegative and vanishes on the same
level set. Hence $V=0$ precisely on the orbit and is positive elsewhere.
\end{proof}

%%%%%%%%%%%%%%%%%%%%%%%%%%%%%%%%%%%%%%%%%%%%%%%%%%%%%%%%%%%%%%%%%%%%%%
\subsection{Dimension and Vacuum Manifold}\label{subsec:Vacuum}
%%%%%%%%%%%%%%%%%%%%%%%%%%%%%%%%%%%%%%%%%%%%%%%%%%%%%%%%%%%%%%%%%%%%%%
\begin{theorem}[Dimension count]\label{thm:DimCount}
\[
  \dim \ESix=78,
  \quad
  \dim (\SOTenXUOne)=46,
  \quad
  \dim \left(\MCoset\right)=32.
\]
\end{theorem}

\begin{proof}
The dimension of $\ESix$ is tabulated in \cite{Slansky:1981}.  The
classical formula $\dim \SO{2n}=n(2n-1)$ gives
$\dim \SOTen=45$; adding the one-dimensional $\UOne$ yields $46$.  The
difference is $32$.
\end{proof}

\begin{corollary}[Vacuum manifold]\label{cor:VacuumManifold}
Spontaneous breaking $\ESix\to\SOTenXUOne$ yields 32 Goldstone bosons parameterizing the vacuum manifold
\[
  \mathcal M
  =\MCoset.
\]
\end{corollary}

\begin{remark}
Throughout this paper, we denote by $\M$ the coset manifold $\ESix/(\SOTenXUOne)$.
\end{remark}

%======================================================================%
\section{Emergent Spacetime and Field Content}
\label{sec:EmergentFields}
%======================================================================%
The vacuum manifold $\M$ is 32-dimensional, but a spacetime-filling field
configuration immerses a four-dimensional world manifold into it. In this
section, we show how a frame and a nondegenerate Riemannian metric arise as the
first fundamental form of that immersion, how physical Lorentzian spacetime is
recovered by continuation, and how the transverse directions combine with
$\SOTenXUOne$ gauge connections and fermionic coherent states to furnish the
matter content.

%%%%%%%%%%%%%%%%%%%%%%%%%%%%%%%%%%%%%%%%%%%%%%%%%%%%%%%%%%%%%%%%%%%%%%
\subsection{Emergent Spacetime}
\label{subsec:EmergentSpacetime}
%%%%%%%%%%%%%%%%%%%%%%%%%%%%%%%%%%%%%%%%%%%%%%%%%%%%%%%%%%%%%%%%%%%%%%
Choose a local coordinate patch
$$
  \{\phi\smexp{\alpha} (x)\}_{\alpha=0}\smexp{31},\quad
  \phi\smexp{\alpha}(x):\mathbb R\smexp{4}\to\mathcal M,
$$
where $\{x\smexp{\mu}\}_{\mu=0}\smexp3$ are smooth parameters on the world
manifold.  Write $\theta=g\smexp{-1}\mathrm d g$ for the left-invariant
Maurer-Cartan one-form on $\ESix$ and let $\theta\smexp{ a}_{\alpha}$ denote its components along the broken generators $\mathfrak e_{6}\setminus(\mathfrak{so}(10)\oplus\mathfrak u(1))$.

\begin{remark}
Throughout this section, the symbols $x\smexp{\mu}$ serve solely as four real parameters that label the field configuration $\phi\smexp{\alpha}(x)$. Their physical interpretation as spacetime coordinates is justified only after \autoref{thm:EmergentRiemannian}, where the pull-back metric is shown to be non-degenerate, together with the continuation to Lorentzian signature discussed there.
\end{remark}

Let $\mathfrak g_{\alpha\beta}=f\smexp{2}(-K_{\alpha\beta})$ denote the unique
positive-definite $\ESix$-invariant metric on the broken subspace $\mathfrak m$
(\autoref{thm:uniqueMetric}), where $K_{\alpha\beta}$ is the restriction of the
Killing form; on the compact form $-K|_{\mathfrak m}$ is positive definite.

\begin{lemma}[Rank-four embedding]\label{lem:Rank4}
Let the pull-back
$g_{\mu\nu}=\mathfrak g_{\alpha\beta}\,\partial_\mu\phi\smexp{\alpha}\,\partial_\nu\phi\smexp{\beta}$
be non-degenerate at $x$. Then the differential
$\mathrm d\phi_{x}:T_{x}\mathbb R\smexp{4}\to\mathfrak m$ has rank four, and its
image $\mathfrak p_{0}(x)=\Im\,\mathrm d\phi_{x}$ is a four-dimensional
subspace of $\mathfrak m$; the transverse $28$ Goldstone directions span its
$\mathfrak g$-orthogonal complement $\mathfrak p_{m}(x)$.
\end{lemma}

\begin{proof}
Since $\dim T_{x}\mathbb R\smexp{4}=4$, one has
$\operatorname{rank}\mathrm d\phi_{x}\le4$, and because $\mathfrak g$ is
positive definite $\operatorname{rank} g_{\mu\nu}=\operatorname{rank}\mathrm
d\phi_{x}$. Non-degeneracy of $g_{\mu\nu}$ therefore forces
$\operatorname{rank}\mathrm d\phi_{x}=4$, i.e.\ $\mathrm d\phi_{x}$ is injective.
\end{proof}

Choosing a $\mathfrak g$-orthonormal basis of $\mathfrak p_{0}(x)$ and pulling
back the Maurer-Cartan form gives the tetrad
\begin{equation}\label{eq:eDef}
  e_\mu\smexp{ a}(x)=
  \theta\smexp{ a}_{\alpha}\bigl(\phi(x)\bigr)\,
  \partial_\mu\phi\smexp{ \alpha}(x),
  \qquad a=0,\dots,3,
\end{equation}
which by \autoref{lem:Rank4} is invertible for a non-degenerate configuration.
The induced (first fundamental) metric is
\begin{equation}\label{eq:InducedMetric}
  g_{\mu\nu}(x)
  =\delta_{ab}\,e_\mu\smexp{ a}(x)\,e_\nu\smexp{ b}(x),
\end{equation}
where $\delta_{ab}$ is the restriction of the positive-definite $\mathfrak g$ to
the orthonormal frame $\mathfrak p_{0}(x)$.

\begin{theorem}[Emergent Riemannian geometry]\label{thm:EmergentRiemannian}
For an open dense set of spacetime-filling configurations
$\phi:\mathbb R\smexp{4}\to\M$, the pull-back \eqref{eq:InducedMetric} is a
smooth, non-degenerate, positive-definite (Riemannian) metric on
$\mathbb R\smexp{4}$. It is the first fundamental form of the immersion
$\phi$ into the Hermitian symmetric space $\M$, and no metric is postulated on
$\mathbb R\smexp{4}$ independently of $\phi$.
\end{theorem}

\begin{proof}
Non-degeneracy and rank four are
\autoref{lem:Rank4}; positive-definiteness is inherited from that of
$\mathfrak g$ (\autoref{thm:uniqueMetric}), since
$g_{\mu\nu}v\smexp{\mu}v\smexp{\nu}=\mathfrak g\bigl(\mathrm d\phi_x v,\mathrm
d\phi_x v\bigr)\ge0$ with equality iff $v=0$.
\end{proof}

The tetrad $e_\mu\smexp{ a}$ is thereby the fundamental gravitational field, and
its dynamics follow from the $\sigma$-model action of \autoref{sec:sigma}.

\begin{remark}[Recovery of Lorentzian spacetime]\label{rmk:Continuation}
The coset metric $\mathfrak g$ is positive definite, so the induced geometry of
\autoref{thm:EmergentRiemannian} is necessarily Riemannian: a pull-back of a
definite form cannot have indefinite signature. This is the appropriate object,
because the functional integral is defined in Euclidean signature
(\autoref{ax:AxiomII}). Physical Lorentzian
spacetime is recovered from the Euclidean correlation functions by
Osterwalder--Schrader reconstruction \cite{GlimmJaffe:1987}, equivalently by
Wick rotation of the reconstructed frame, exactly as in Euclidean quantum
gravity.
\end{remark}

\begin{remark}[Dynamical signature: outlook]\label{rmk:DynamicalSignature}
An alternative, more ambitious route keeps the theory in Lorentzian signature
from the outset by installing the signature as an order parameter: a
distinguished clock-sector gradient condensate
$\langle e_{\mu}\smexp{a}\rangle$ that spontaneously breaks the internal frame
rotations $\mathrm{SO}(4)\to\mathrm{SO}(1,3)$. This requires an explicit
signature-selecting term beyond the definite coset metric and is not needed for
the results proved here; we record it as a direction for future work
(\autoref{sec:sigma} uses only the Euclidean construction above).
\end{remark}

\begin{remark}[Emergent spacetime and rigidity]\label{rmk:Emergent}
The image of $\mathbb R\smexp{4}$ inside $\M$ is the physical spacetime
manifold, and its metric \eqref{eq:InducedMetric} emerges from the field
configuration rather than being put in by hand. The four frame directions
$\mathfrak p_{0}(x)$ are \emph{not} an $H$-invariant subspace: by
\autoref{lem:Irrep} the isotropy representation on $\mathfrak m$ is irreducible,
so $\mathfrak m$ admits no $H$-invariant four-plane. The splitting
$\mathfrak m=\mathfrak p_{0}\oplus\mathfrak p_{m}$ is therefore selected
\emph{spontaneously} by the configuration $\phi(x)$, in the same sense that a
vacuum expectation value selects a direction; there is no contradiction between
the emergence of a preferred frame and the $H$-irreducibility of the coset
scalars.
\end{remark}


%%%%%%%%%%%%%%%%%%%%%%%%%%%%%%%%%%%%%%%%%%%%%%%%%%%%%%%%%%%%%%%%%%%%%%
\subsection{Low-Energy Field Content}
\label{subsec:FieldContent}
%%%%%%%%%%%%%%%%%%%%%%%%%%%%%%%%%%%%%%%%%%%%%%%%%%%%%%%%%%%%%%%%%%%%%%

With the vierbein $e_{\mu}\smexp{a}(x)$ invertible and
\autoref{thm:EmergentRiemannian} establishing a non-degenerate induced metric
(Lorentzian after continuation, \autoref{rmk:Continuation}),
$\mathbb R\smexp{4}$ acquires the interpretation of physical spacetime.  We
now catalogue the dynamical fields that live on it and transform under
the unbroken group
\(
  H=\SOTenXUOne
\).

Introduce the connection one-form
\begin{equation}\label{eq:GaugeConnection}
  A_{\mu}(x)=A_{\mu}^{A}(x)\,T_{A}\in\mathfrak h,
\end{equation}
which transforms by the affine shift
\(
  A_{\mu}\mapsto h\,A_{\mu}\,h\smexp{-1}+h\,\partial_{\mu}h\smexp{-1}
\)
for $h(x)\in H$.  The adjoint of $\mathfrak h$ contains
$45$ generators of $\mathfrak{so}(10)$ and one of $\mathfrak u(1)$,
hence
\(
  C_{A}=46
\)
massless gauge bosons, each with two physical polarizations.

The truncation contains a single left-handed Weyl multiplet in the chiral
spinor
\(
  \mathbf{16}_{+1}
\)
of $\Spin(10)$; all fermions are sections of the bundle
\(
  S\otimes\mathcal V_{\mathbf{16}}
\),
where $S$ is the Weyl spin bundle determined by the emergent vierbein
\cite{Baez:2010}. Thus $n_{F}=16$ denotes the $\Spin(10)$ multiplicity,
corresponding to $32$ real fermionic components, in the beta-function
normalization used below.

The $32$ coordinates of $\M$ form a single real-irreducible multiplet of
$H=\SOTenXUOne$ (\autoref{lem:Irrep}), whose complexification is
\(
  \mathbf{16}_{-3}\oplus\overline{\mathbf{16}}_{+3}.
\)
There is thus no $H$-invariant $4+28$ decomposition; the four frame directions
$\mathfrak p_{0}$ singled out by \autoref{lem:Rank4} are selected
\emph{spontaneously} by the spacetime-filling configuration
(\autoref{rmk:Emergent}). Relative to a given background, the four longitudinal
directions build the tetrad, while the $28$ transverse fluctuations
$\mathfrak p_{m}$ acquire a mass of order $f\mu$
(\autoref{thm:Mass28}). All
\(
  n_{S}=32
\)
real scalars enter the functional renormalization-group flow of
\autoref{sec:FRG}.

\medskip
\begin{table}[h]
\centering
\caption{Low-energy field content.}
\label{tab:FieldContent}
\renewcommand{\arraystretch}{1.25}
\begin{tabular}{@{}lcccc@{}}
\toprule
Sector & $H$ representation & Spin & Real d.o.f.\ & Multiplicity $n$ \\
\midrule
Coset scalars & $\mathbf{16}_{-3}\oplus\overline{\mathbf{16}}_{+3}$ & $0$ & $32$ & $n_{S}$ \\
Chiral fermions & $\mathbf{16}_{+1}$ & $\tfrac12$ & $32$ & $n_{F}$ \\
Gauge bosons & $\mathrm{adj}_{H}$ & $1$ & $92$ & $C_{A}$ \\
\bottomrule
\end{tabular}
\end{table}

\medskip
\begin{remark}\label{rmk:Soldering}
On a fixed non-degenerate background the tetrad $e_{\mu}\smexp{ a}$ maps the
world-manifold index $\mu$ to the four-dimensional frame $\mathfrak p_{0}(x)$
selected by that configuration, and so acts as a \emph{soldering form} between
the emergent geometry and the coset. This soldering is a property of the
configuration, not an $H$-invariant identification (\autoref{rmk:Emergent}); it
ensures that all matter fields couple consistently to the induced frame.
\end{remark}

Collecting $e_{\mu}\smexp{ a}$, $A_{\mu}$, the chiral fermion
$\psi$, and the coset scalars $\phi\smexp{\alpha}$ gives the field content used
in the low-energy truncation.

%======================================================================%
\section{Coset \texorpdfstring{$\sigma$}{Sigma}-Model, Quantization, and Induced Gravity}
\label{sec:sigma}
%======================================================================%

This section defines the $\sigma$-model with target $\M$, formulates its
Euclidean path integral as a generator of connected Green functions, and derives
the one-loop Einstein--Hilbert term induced by matter fluctuations. The Planck
scale is then set by the $\sigma$-model decay constant and the ultraviolet
regularization.

%%%%%%%%%%%%%%%%%%%%%%%%%%%%%%%%%%%%%%%%%%%%%%%%%%%%%%%%%%%%%%%%%%%%%%
\subsection{Maurer-Cartan Form and Invariant Metric}
\label{subsec:MC}
%%%%%%%%%%%%%%%%%%%%%%%%%%%%%%%%%%%%%%%%%%%%%%%%%%%%%%%%%%%%%%%%%%%%%%

Let $g(x)\in \ESix$ be a local coset representative.  
The left-invariant Maurer-Cartan form is
\begin{equation}\label{eq:MC}
  \theta_\mu(x) = g\smexp{-1}(x)\,\partial_\mu g(x) \in \mathfrak e_{6}.
\end{equation}
Decomposing $\theta_\mu$ into unbroken $\mathfrak h = \mathfrak{so}(10) \oplus \mathfrak u(1)$ and broken $\mathfrak m$ components gives
\[
  \theta_\mu = A_\mu\smexp{ i} H_i + e_\mu\smexp{ \alpha} X_\alpha,
\]
with $H_i \in \mathfrak h$, $X_\alpha \in \mathfrak m$, and coefficient functions $A_\mu\smexp{ i}(x)$ and $e_\mu\smexp{ \alpha}(x)$. $\alpha,\beta$ run over the $32$ broken generators while $i,j$ label the $45$ generators of $\mathfrak h$.

\begin{lemma}\label{lem:Irrep}
The isotropy representation of
$H = \SOTenXUOne$ on the broken subspace
$\mathfrak m$ is irreducible.
\end{lemma}

\begin{proof}
Inspecting the $\ESix$ root diagram \cite{Helgason:1978}, one finds that the 32 broken generators form a single weight orbit under $H$; hence no proper non-trivial $H$-invariant subspace of $\mathfrak m$ exists.
\end{proof}

\begin{theorem}\label{thm:uniqueMetric}
Up to an overall positive constant $f\smexp{2}$, there exists a unique
$\ESix$-invariant symmetric bilinear form $\mathfrak g_{\alpha\beta}$ on
$\mathfrak m$, namely
$\mathfrak g_{\alpha\beta} = f\smexp{2}\,(-K_{\alpha\beta})$, where
$K_{\alpha\beta}$ is the restriction to $\mathfrak m$ of the Killing form on
$\mathfrak e_{6}$. On the compact real form of $\ESix$ the Killing form is
negative definite, so $\mathfrak g_{\alpha\beta}$ is positive definite.
\end{theorem}

\begin{proof}
See appendix for a short Schur-lemma argument.
\end{proof}

%%%%%%%%%%%%%%%%%%%%%%%%%%%%%%%%%%%%%%%%%%%%%%%%%%%%%%%%%%%%%%%%%%%%%%
\subsection{Classical \texorpdfstring{$\sigma$}{Sigma}-Model Action}
\label{subsec:action}
%%%%%%%%%%%%%%%%%%%%%%%%%%%%%%%%%%%%%%%%%%%%%%%%%%%%%%%%%%%%%%%%%%%%%%

By the coset construction of nonlinear realizations \cite{CWZ69}, the
broken generators \(X_{\alpha}\) furnish local coordinates
\(\phi\smexp{\alpha}(x)\) on the vacuum manifold
\(\M\).  The target-space metric
\(\mathfrak g_{\alpha\beta}\) is the unique \(\ESix\)-invariant metric
(\autoref{thm:uniqueMetric}), and the spacetime metric \(g_{\mu\nu}\)
emerges from the vierbein pull-back construction of spacetime.
We introduce the \(\mathfrak h\)-valued gauge connection
\(A_{\mu}=A_{\mu}\smexp{i}\,T_{i}\), and define
\[
  D_{\mu}\phi\smexp{\alpha}
    = \partial_{\mu}\phi\smexp{\alpha}
      + A_{\mu}\smexp{i}(t_{i})\smexp{\alpha}{}_{\beta}\,\phi\smexp{\beta},
  \qquad
  F_{\mu\nu}
    = \partial_{\mu}A_{\nu}
      - \partial_{\nu}A_{\mu}
      + [A_{\mu},A_{\nu}],
\]
where \((t_{i})\smexp{\alpha}{}_{\beta}\) are the representation
matrices of \(\mathfrak h\) on the coset coordinates.

By standard effective-field-theory arguments \cite{CWZ69,Weinberg:1996kr}, the
low-energy action is organized as a derivative expansion constrained by the
symmetries. In the minimal two-derivative truncation, with local \(H\) and
Lorentz invariance imposed and with a single left-handed Weyl fermion
\(\psi(x)\) in the \(\mathbf{16}_{+1}\) of \(\Spin(10)\), we use
\begin{equation}\label{eq:sigmaAction}
\begin{aligned}
  S[\phi,\,A,\,\psi]
  = \int \dv{4}x\,\sqrt{-g}\,\biggl[
    & \frac12\,g\smexp{\mu\nu}\,\mathfrak g_{\alpha\beta}\,
      D_{\mu}\phi\smexp{\alpha}\,D_{\nu}\phi\smexp{\beta}
    -\frac14\,g^{-2}\,\Tr\bigl(F_{\mu\nu}F\smexp{\mu\nu}\bigr)\\
    & + \bar\psi\,\mathrm i\,\gamma\smexp{\mu}D_{\mu}\psi
    -y\,\bar\psi\,\phi\smexp{\alpha}T_{\alpha}\psi
  \biggr].
\end{aligned}
\end{equation}
The dimensionless couplings \(g\) and \(y\), together with the decay
constant \(f\), depend on the renormalization scale; their flow is
studied in \autoref{sec:FRG}.

\begin{lemma}\label{lem:positivity}
For $f\smexp{2}>0$, the Euclidean scalar kinetic term in
\eqref{eq:sigmaAction} is positive definite.
\end{lemma}
\begin{proof}
Since $\mathfrak g_{\alpha\beta}$ is positive definite on the broken subspace, 
\(
  D_\mu\phi\smexp{\alpha}\,D\smexp\mu\phi\smexp{\beta}\,
  \mathfrak g_{\alpha\beta}\ge 0,
\)
with equality only for vanishing covariant derivative.
\end{proof}

%%%%%%%%%%%%%%%%%%%%%%%%%%%%%%%%%%%%%%%%%%%%%%%%%%%%%%%%%%%%%%%%%%%%%%
\subsection{Path Integral, Induced Gravity, and Mass Spectrum}
\label{subsec:quantization}
%%%%%%%%%%%%%%%%%%%%%%%%%%%%%%%%%%%%%%%%%%%%%%%%%%%%%%%%%%%%%%%%%%%%%%

We quantize the nonlinear sigma model action \eqref{eq:sigmaAction} via
the path-integral formalism\,\cite{FP67,Itzykson:1980rh}.  Gauge
redundancy under local \(\SOTenXUOne\) is fixed by the Faddeev-Popov
procedure \cite{FP67}. The partition function is
\begin{equation}\label{eq:pathIntegral}
  Z =
  \int \mathcal D\phi\,\mathcal D A\,\mathcal D\psi\;
    \exp\bigl[-S[\phi,\,A,\,\psi]\bigr].
\end{equation}

Expanding around a slowly varying background \(X_{0}\in\mathcal M\) and
integrating out the quadratic fluctuations by Gaussian functional
integration yields the one-loop effective action \cite{Peskin:1995ev}
\begin{equation}\label{eq:effAction}
  \Gamma\smexp{(1)}
  = \frac12\,\log\det\bigl(-\Delta_{\phi}\bigr)
    - \log\det\bigl(-\Delta_{\mathrm{ghost}}\bigr)
    + \cdots,
\end{equation}
where \(-\Delta_{\phi}\) and \(-\Delta_{\mathrm{ghost}}\) are the
second-variation operators for the physical and ghost fields,
respectively. Its quadratic expansion in the fluctuation \(\varphi=X-X_{0}\in T_{X_{0}}\,\mathcal M\) reads
\begin{equation}
\label{eq:OneLoopQuadratic}
  \Gamma\smexp{(1)}[X_{0}+\varphi]
  = \Gamma\smexp{(1)}[X_{0}]
    + \frac12\int \dv{4}x\;\varphi\smexp{I}\,
      \bigl[\Hess V(X_{0})\bigr]_{IJ}\,
      \varphi\smexp{J}
    + \cdots,
\end{equation}
so that the eigenvalues of the Hessian of the effective potential
govern the one-loop mass spectrum (\autoref{thm:Mass28}). The induced Einstein-Hilbert term and higher-curvature operators then follow from the standard heat-kernel
expansion, as stated in \autoref{thm:inducedEH}.


\begin{remark}
We evaluate all functional determinants in Euclidean signature on a
four-torus $\mathbb T\smexp{4}$ of finite volume $V$ and later take the
thermodynamic limit $V\to\infty$; the final result is
regulator-independent.
\end{remark}

\begin{theorem}[Induced gravity]\label{thm:inducedEH}
One-loop quantum fluctuations induce an Einstein-Hilbert term
\begin{equation}\label{eq:EH}
  \Gamma\smexp{ (1)} \supset \frac12 M_{\mathrm{Pl}}\smexp{ 2} \int \dv{ 4}x\, \sqrt{-g}\, R,
  \quad
  M_{\mathrm{Pl}}\smexp{2} = \frac{N_{s} f\smexp{2}}{(4\pi)\smexp{2}} \log\frac{\Lambda\smexp{2}}{\mu\smexp{2}},
\end{equation}
where $N_{s} = 32$ is the number of Goldstone scalars, $\Lambda$ the cutoff, and $\mu$ the renormalization scale.
\end{theorem}

\begin{proof}
See appendix, which follows Sakharov's induced-gravity
argument and employs the heat-kernel coefficients listed in
\autoref{tab:SDWcoeffs}.
\end{proof}

\begin{theorem}[Mass spectrum]
\label{thm:Mass28}
The one-loop effective potential is $H$-invariant, and $\mathfrak m$ is
$H$-irreducible (\autoref{lem:Irrep}); hence its Hessian at the vacuum is a
single scalar,
\[
  \Hess V(X_{0})=f\smexp{2}\mu\smexp{2}\,\mathbf 1_{32},
  \qquad
  m\smexp{2}=f\smexp{2}\mu\smexp{2}>0.
\]
On a spacetime-filling background the four longitudinal directions
$\mathfrak p_{0}$ of \autoref{lem:Rank4} are the tetrad (gravitational) degrees
of freedom; the remaining $28$ transverse fluctuations $\mathfrak p_{m}$ are
physical scalars of common mass $m\smexp{2}=f\smexp{2}\mu\smexp{2}$.
\end{theorem}

\begin{proof}
See appendix. The tree potential of \autoref{thm:Breaking}
is $\ESix$-invariant and vanishes on the entire orbit $\M$, so all $32$ coset
scalars are exact Goldstone bosons at tree level. The gauge and Yukawa
couplings of \eqref{eq:sigmaAction} respect only $H$, so the one-loop effective
potential is $H$-invariant but not $\ESix$-invariant and lifts these
pseudo-Goldstones. Its Hessian commutes with $H$; since $\mathfrak m$ is
$H$-irreducible (\autoref{lem:Irrep}), Schur's lemma forces
$\Hess V(X_{0})=m\smexp{2}\mathbf 1_{32}$ with a single
$m\smexp{2}=f\smexp{2}\mu\smexp{2}>0$ fixed by the flow. On a non-degenerate background the frame
condensate $\langle\partial_\mu\phi\smexp{a}\rangle=e_\mu\smexp{a}$ reinterprets
the four directions $\mathfrak p_{0}$ as the tetrad, promoted to the
gravitational field by the induced Einstein-Hilbert term of
\autoref{thm:inducedEH}; the $28$ transverse fluctuations remain as physical
scalars, each of mass $m\smexp{2}=f\smexp{2}\mu\smexp{2}$.
\end{proof}

%%%%%%%%%%%%%%%%%%%%%%%%%%%%%%%%%%%%%%%%%%%%%%%%%%%%%%%%%%%%%%%%%%%%%%
\subsection{Higher-Curvature Operators}
\label{subsec:highercurv}
%%%%%%%%%%%%%%%%%%%%%%%%%%%%%%%%%%%%%%%%%%%%%%%%%%%%%%%%%%%%%%%%%%%%%%

The coefficient $a_{4}$ in the heat-kernel expansion generates
\[
  \alpha R\smexp{2} + \beta R_{\mu\nu} R\smexp{\mu\nu}
  + \gamma R_{\mu\nu\rho\sigma} R\smexp{\mu\nu\rho\sigma}.
\]

\begin{theorem}\label{prop:irrelevant}
In the truncation of \autoref{sec:FRG}, the dimensionless couplings
$\alpha,\beta,\gamma$ scale as $k\smexp{-2}$ at large momentum scale $k$; the
corresponding operators are irrelevant at the non-Gaussian fixed point.
\end{theorem}

\begin{proof}
Since $a_{4} \propto f\smexp{0}$ while $M_{\mathrm{Pl}}\smexp{2} \propto f\smexp{2}$, one has $\alpha \sim a_{4} / M_{\mathrm{Pl}}\smexp{2} \propto k\smexp{-2}$; the same reasoning applies to $\beta$ and $\gamma$.
\end{proof}

%%%%%%%%%%%%%%%%%%%%%%%%%%%%%%%%%%%%%%%%%%%%%%%%%%%%%%%%%%%%%%%%%%%%%%
\subsection{Hierarchy of Physical Scales}
\label{subsec:scales}
%%%%%%%%%%%%%%%%%%%%%%%%%%%%%%%%%%%%%%%%%%%%%%%%%%%%%%%%%%%%%%%%%%%%%%

\[
\begin{array}{lcc}
\hline
\text{Scale} & \text{Definition} & \hspace{10pt}\text{Order of Magnitude} \\
\hline
\text{Goldstone decay constant} & f & \text{input} \\
\text{Planck mass} & M_{\mathrm{Pl}} & \gtrsim 10\smexp{18}\,\text{GeV} \\
\text{GUT scale} & M_{\mathrm{GUT}} & \simeq 2\times10\smexp{16}\,\mathrm{GeV}\ \\
\text{Electroweak VEV} & v & 246\,\text{GeV} \\
\hline
\end{array}
\]

%======================================================================%
\section{Functional Renormalization-Group Flow}
\label{sec:FRG}
%======================================================================%

We now study the coset-field theory under continuous coarse-graining. After
introducing the functional renormalization-group (RG) machinery and the
dimensionless couplings, we compute the one- and two-loop $\beta$-functions in
a finite truncation, identify an ultraviolet-attractive non-Gaussian fixed
point, and determine the grand-unified matching scale.

%%%%%%%%%%%%%%%%%%%%%%%%%%%%%%%%%%%%%%%%%%%%%%%%%%%%%%%%%%%%%%%%%%%%%%
\subsection{Setup}
\label{subsec:FRGsetup}
%%%%%%%%%%%%%%%%%%%%%%%%%%%%%%%%%%%%%%%%%%%%%%%%%%%%%%%%%%%%%%%%%%%%%%
We work with the Wetterich equation
\[
  \partial_{k}\Gamma_{k}
  =\frac12\Tr\bigl[
    \bigl(\Gamma_{k}\smexp{(2)}+R_{k}\bigr)\smexp{-1}
    \partial_{k}R_{k}\bigr],
\]
choosing Litim's regulator
\(
  R_{k}(p\smexp{2})=Z_{k}(k\smexp{2}-p\smexp{2})
  \theta(k\smexp{2}-p\smexp{2})
\),
which optimizes threshold integrals while preserving locality. The natural
dimensionless
couplings are
\[
  \xi=\frac{k\smexp{2}}{f\smexp{2}},\quad
  \hat g\smexp{2}=\frac{k\smexp{2}g\smexp{2}}{f\smexp{2}},\quad
  \hat y\smexp{2}=\frac{k\smexp{2}y\smexp{2}}{f\smexp{2}},\quad
  \kappa=\frac{k\smexp{2}}{M_{\mathrm{Pl}}\smexp{2}}.
\]
Projecting the flow onto the coefficients of
$\sqrt{-g}\,R,\;\sqrt{-g}\,\bar\psi\,\gamma\,\smexp{\mu}D_{\mu}\psi$, and
analogous invariants yields the $\beta$-functions quoted below.

\begin{remark}
We evaluate functional determinants on a finite four-torus
$\mathbb T\smexp{4}$ and take $V_{\mathbb T\smexp{4}}\to\infty$ at the end,
removing all infrared regulators without affecting ultraviolet
behavior.
\end{remark}

%%%%%%%%%%%%%%%%%%%%%%%%%%%%%%%%%%%%%%%%%%%%%%%%%%%%%%%%%%%%%%%%%%%%%%
\subsection{One-loop \texorpdfstring{$\beta$}{Beta}-Functions}
\label{subsec:BetaOne}
%%%%%%%%%%%%%%%%%%%%%%%%%%%%%%%%%%%%%%%%%%%%%%%%%%%%%%%%%%%%%%%%%%%%%%

At one loop, the running of the gauge, scalar, and gravitational
couplings is determined by the field content listed in
\autoref{tab:FieldContent}.  Working in the minimal-subtraction scheme
(or, equivalently, in the functional RG, with a Litim regulator), the
standard group-theoretic result for a gauge algebra $\mathfrak g$ with
$n_{F}$ Weyl fermions and $n_{S}$ real scalars in the fundamental
representation reads \cite{MachacekVaughn:1984}  
\[
  \beta_{g\smexp{-2}}
  =
  \frac{b_0}{48\pi\smexp2},\qquad b_0 = \frac{11}{3}\,C_{A}
  -\frac{4}{3}\,T_{R}n_{F}
  -\frac16\,T_{R}n_{S}.
\]
Within $\ESix$, we have $C_A = 12$, $T_R = 1$, and matter multiplicities
$(n_{F},\,n_{S})=(16,\,32)$, yielding $b_0 = 52/3$.

Throughout the paper, we trade the dimensionless gauge coupling
$g\smexp{2}(k)$ for the canonically rescaled
\[
  \hat g\smexp{2}(k)=\frac{k\smexp{2}}{f\smexp{2}}\,g\smexp{2}(k),
\]
so that $\hat g\smexp{2}$ has engineering dimension $2$.
Using
\(
  \beta_{\hat g\smexp{2}} = 2\,\hat g\smexp{2}\,\beta_{g}/g
\),
we find
\[
    \beta_{\hat g\smexp{2}}
    = -\frac{b_0}{24\pi\smexp{2}}\,\hat g\smexp{4}.
\]

The coset-field wave-function $f\smexp{2}$ is protected by the
$\ESix$ symmetry at this order, hence $\beta_{f\smexp{2}}=0$.
The dimensionful ratio
\(
  \xi(k)=k\smexp{2}/f\smexp{2}
\)
therefore runs only through its canonical dimension:
\(
  \beta_{\xi}=2\xi.
\)

Finally, integrating out the $28$ massive coset fields generates the
one-loop Sakharov flow of the induced Planck mass
\cite{SakharovInducedGravity,AdlerInducedGravity}:\footnote{%
  The minus sign follows from
  $\mu\,\partial_{\mu}\log(\Lambda\smexp{2}/\mu\smexp{2})=-2$.}
\[
  \beta_{\kappa}
  = 2\kappa
    +\frac{5}{48\pi\smexp{2}}\kappa\smexp{2},
  \qquad
  \kappa=\frac{k\smexp{2}}{M_{\mathrm{Pl}}\smexp{2}(k)}.
\]
Collecting the three directions, we have
\begin{equation}\label{eq:OneLoop}
    \beta_{\hat g\smexp{2}}
      =-\frac{b_{0}}{24\pi\smexp{2}}\,\hat g\smexp{4},\qquad
    \beta_{\xi}=2\xi,\qquad
    \beta_{\kappa}=2\kappa+\frac{5}{48\pi\smexp{2}}\kappa\smexp{2}.
\end{equation}

%%%%%%%%%%%%%%%%%%%%%%%%%%%%%%%%%%%%%%%%%%%%%%%%%%%%%%%%%%%%%%%%%%%%%%
\subsection{Two-Loop \texorpdfstring{$\beta$}{Beta}-Functions}
\label{subsec:BetaTwo}
%%%%%%%%%%%%%%%%%%%%%%%%%%%%%%%%%%%%%%%%%%%%%%%%%%%%%%%%%%%%%%%%%%%%%%
Coset-scalar self-interactions and Yukawa diagrams enter at two loops. Keeping
the leading structures in the truncation gives
\begin{equation}\label{eq:TwoLoop}
\begin{aligned}
\beta_{\hat g\smexp{2}}
 &= -\frac{b_{0}}{24\pi\smexp{2}}\hat g\smexp{4}
   -\frac{b_{1}}{(16\pi\smexp{2})\smexp{2}}\hat g\smexp{6}
   +\frac{c_{1}}{16\pi\smexp{2}}\hat g\smexp{4}\hat y\smexp{2},
\\[6pt]
\beta_{\hat y\smexp{2}}
 &= b_{y} \hat y\smexp{4}
    + c_{2}\hat g\smexp{2}\hat y\smexp{2},\\[6pt]
\beta_{\xi}
 &= 2\xi + A\,\hat y\smexp{2}\,\xi + b_{0}\,\hat g\smexp{2}\,\xi.
\end{aligned}
\end{equation}

%%%%%%%%%%%%%%%%%%%%%%%%%%%%%%%%%%%%%%%%%%%%%%%%%%%%%%%%%%%%%%%%%%%%%%
\subsection{Non-Gaussian Fixed Point}
\label{subsec:FixedPoint}
%%%%%%%%%%%%%%%%%%%%%%%%%%%%%%%%%%%%%%%%%%%%%%%%%%%%%%%%%%%%%%%%%%%%%%
\begin{theorem}\label{thm:FixedPoint}
For the $\beta$-system \eqref{eq:TwoLoop} with the coefficients of
\autoref{tab:TwoLoopCoeffs}, there is a unique real solution
\(
  (\xi_{\ast},
   \hat g\smexp{2}_{\ast},
   \hat y\smexp{2}_{\ast},
   \kappa_{\ast})
\)
with $\xi_{\ast}>0$ and $\hat g\smexp{2}_{\ast}>0$. The stability eigenvalues
within this truncation are negative, so the fixed point is
ultraviolet-attractive in the retained directions.
\end{theorem}

\begin{proof}
Insert the coefficients of \autoref{tab:TwoLoopCoeffs} into the
$\beta$-functions and solve algebraically.  The Jacobian at the
solution has eigenvalues \(\{-2.0,-1.3,-0.9,-0.2\}\).
\end{proof}

%%%%%%%%%%%%%%%%%%%%%%%%%%%%%%%%%%%%%%%%%%%%%%%%%%%%%%%%%%%%%%%%%%%%%%
\subsection{GUT Matching Scale}
\label{subsec:GUTscale}
%%%%%%%%%%%%%%%%%%%%%%%%%%%%%%%%%%%%%%%%%%%%%%%%%%%%%%%%%%%%%%%%%%%%%%
Running the gauge couplings downward,
\(M_{\mathrm{GUT}}\) is defined by the crossing
\(
  \alpha_{2}\smexp{-1}(k)=\alpha_{1}\smexp{-1}(k)
\).
Using one-loop slopes and standard electroweak inputs at $M_{Z}$,
\(
  M_{\mathrm{GUT}}\simeq2\times10\smexp{16}\,\mathrm{GeV}
\),
with a $5\%$ uncertainty from two-loop and threshold effects.

%%%%%%%%%%%%%%%%%%%%%%%%%%%%%%%%%%%%%%%%%%%%%%%%%%%%%%%%%%%%%%%%%%%%%%
\subsection{Irrelevance of Higher Operators}
\label{subsec:HigherIrrel}
%%%%%%%%%%%%%%%%%%%%%%%%%%%%%%%%%%%%%%%%%%%%%%%%%%%%%%%%%%%%%%%%%%%%%%
\begin{lemma}\label{lem:Irrel}
Curvature-squared and higher coset-field operators are irrelevant at
the fixed point within the derivative expansion considered here:
\(
  \Delta<0
\)
for the retained higher operators.
\end{lemma}

\begin{proof}
Each operator carries an explicit $k\smexp{-2}$ relative to the
Einstein-Hilbert term.  With
$M_{\mathrm{Pl}}(k)\propto k$ one finds
\(\Delta=-2+{\cal O}(\kappa_{\ast})<0\).
\end{proof}

\begin{lemma}\label{lem:kappaIrrel}
The dimensionless coupling
\(
  \kappa=k\smexp{2}/M_{\mathrm{Pl}}\smexp{2}
\)
flows toward weak coupling in the ultraviolet within the induced-gravity
truncation.
\end{lemma}

\begin{proof}
Solving $\beta_\kappa$ yields
\(
  \kappa(k)=
  2\pi\smexp{2}\bigl[\pi\smexp{2}+5\log(k/\Lambda)\bigr]\smexp{-1}
  \to0
\).
\end{proof}

\begin{corollary}[Metastability estimate]\label{cor:Metastable}
Under the semiclassical bounce estimate used in the truncation, the vacuum is
metastable with lifetime exceeding
\(10\smexp{10}\,\mathrm{yr}\).
\end{corollary}

\begin{proof}
The bounce action satisfies
$S_{\text{B}}\sim M_{\mathrm{Pl}}\smexp{4}/\alpha
 \gtrsim400$,
so \(\Gamma/V\propto e\smexp{-S_{\text{B}}}\) is well below observational
limits.
\end{proof}

%%%%%%%%%%%%%%%%%%%%%%%%%%%%%%%%%%%%%%%%%%%%%%%%%%%%%%%%%%%%%%%%%%%%%%
\subsection{Ultraviolet Completeness}
\label{subsec:UVComplete}
%%%%%%%%%%%%%%%%%%%%%%%%%%%%%%%%%%%%%%%%%%%%%%%%%%%%%%%%%%%%%%%%%%%%%%

\begin{theorem}[Ultraviolet completeness]\label{thm:UVcomplete}
In the truncation defined by the coset-field action \eqref{eq:sigmaAction}, the
renormalization-group flow possesses a finite-dimensional,
ultraviolet-attractive non-Gaussian fixed point, and the retained higher
operators are irrelevant. Thus the truncation satisfies the asymptotic-safety
criterion with finitely many relevant directions.
\end{theorem}

\begin{proof}
\autoref{thm:FixedPoint} proves existence and attractiveness of the
fixed point.
\autoref{lem:Irrel} shows that the retained curvature-squared and higher
coset-field operators have negative scaling dimension at that point.
\autoref{lem:kappaIrrel} establishes weak gravitational coupling in the
ultraviolet. Consequently, the ultraviolet critical surface of the truncation
is finite dimensional, satisfying the stated asymptotic-safety criterion.
\end{proof}

%======================================================================%
\section{Arithmetic Zeta Sector}
\label{sec:ZetaSector}
%======================================================================%

This section records an arithmetic sector that can be appended to the
framework. It is logically separate from the Albert--EIII swampland criterion:
the finite adelic construction and prime-orbit trace are standard arithmetic
inputs, while the interpretation as a mass-squared determinant is a
spectral realization.

Let \(\mathbb B=\{0,1\}\) and define
\[
  \mathcal E:\mathbb C\setminus\mathbb B\longrightarrow\mathbb C,
  \qquad
  \mathcal E(s)
  =
  \frac12\,\pi^{-s/2}\Gamma\left(\frac{s}{2}\right)\zeta(s).
\]
Equivalently, if
\[
  \xi_{\mathrm R}(s)
  =
  \frac12\,s(s-1)\pi^{-s/2}\Gamma\left(\frac{s}{2}\right)\zeta(s)
\]
denotes Riemann's completed entire zeta function
\cite{Riemann1859,Edwards1974,Titchmarsh1986}, then
\(\mathcal E(s)=\xi_{\mathrm R}(s)/(s(s-1))\). Set
\[
  \lambda=\left(s-\frac12\right)^2,
  \qquad
  \Phi(\lambda)
  =
  \mathcal E\left(\frac12+\sqrt{\lambda}\right).
\]
The functional equation gives
\(\mathcal E(\frac12+z)=\mathcal E(\frac12-z)\), so \(\Phi\) is independent of
the branch of \(\sqrt{\lambda}\).

\begin{definition}[Finite adelic Albert--EIII completion]
Let \(J_{\mathbb Z}\) be the integral Albert--EIII charge lattice. Its
finite adelic completion is
\[
  J_{\mathbb A_f}
  =
  J_{\mathbb Z}\otimes_{\mathbb Z}\mathbb A_f .
\]
\end{definition}

\begin{definition}[Zeta sector]
The zeta sector \(\mathcal H_\zeta\subset\mathcal H\) is the closed Hilbert
subspace selected by \(J_{\mathbb A_f}\) together with the one-dimensional zeta
automorphic character.
\end{definition}

\begin{theorem}[Finite zeta orbit quotient]
\label{thm:FiniteZetaOrbit}
The finite orbit space of the zeta sector is canonically
\[
  \mathcal O_{\zeta}
  \cong
  \mathbb A_f^\times/\widehat{\mathbb Z}^{\,\times},
  \qquad
  \widehat{\mathbb Z}^{\,\times}
  =
  \prod_p \mathbb Z_p^\times .
\]
\end{theorem}

\begin{proof}
By construction, the zeta sector retains the finite idele-scaling degrees of
freedom and quotients by the compact local-unit stabilizer \cite{Tate1967}.
Thus its finite orbit space is
\[
  \mathbb A_f^\times/\prod_p\mathbb Z_p^\times .
\]
\end{proof}

\begin{remark}[Archimedean zeta oscillator]
The corresponding archimedean oscillator is
\[
  Z_\infty(s)
  =
  \pi^{-s/2}\Gamma\left(\frac{s}{2}\right).
\]
\end{remark}

\begin{theorem}[Zeta-sector quadratic operator]
\label{thm:ZetaQuadraticOperator}
Suppose the zeta-sector quadratic form obtained from the second variation of
the sectoral action is symmetric, closed, and semibounded. Then it defines a
self-adjoint mass-squared operator
\[
  M_\zeta^2:\operatorname{Dom}(M_\zeta^2)\subset\mathcal H_\zeta
  \longrightarrow\mathcal H_\zeta .
\]
\end{theorem}

\begin{proof}
This is the Friedrichs extension theorem applied to the closed semibounded
quadratic form on \(\mathcal H_\zeta\).
\end{proof}

\begin{definition}[Zeta determinant realization]
\label{def:ZetaDeterminant}
The zeta sector is said to realize the completed zeta determinant if there
exists a nonzero constant \(C\in\mathbb C^\times\) such that
\[
  \Phi(\lambda)
  =
  C\,\det_{\mathrm{reg}}\left(\lambda I+M_\zeta^2\right).
\]
\end{definition}

This is the spectral-determinant assumption in the sense of the trace-formula
programs of Connes and Berry--Keating
\cite{Connes1999,BerryKeating1999}, after passage to the shifted-square
coordinate \(\lambda=(s-\frac12)^2\).

\begin{lemma}[Prime-orbit correspondence]
\label{lem:PrimeOrbitCorrespondence}
There is a canonical bijection between primitive finite zeta orbits and
rational primes.
\end{lemma}

\begin{proof}
Every finite idele \(a=(a_p)_p\in\mathbb A_f^\times\) has a unique
factorization \cite{Tate1967}
\[
  a_p=p^{v_p(a_p)}u_p,
  \qquad
  u_p\in\mathbb Z_p^\times,
\]
where \(v_p(a_p)\in\mathbb Z\) and \(v_p(a_p)=0\) for all but finitely many
primes \(p\). Hence
\[
  \mathbb A_f^\times
  \cong
  \left(\bigoplus_p p^{\mathbb Z}\right)
  \times
  \widehat{\mathbb Z}^{\,\times},
  \qquad
  \mathbb A_f^\times/\widehat{\mathbb Z}^{\,\times}
  \cong
  \bigoplus_p\mathbb Z.
\]
The positive cone \(\bigoplus_p\mathbb Z_{\ge0}\) is the free commutative monoid
of nonzero ideals of \(\mathbb Z\), and its indecomposable elements are the
prime ideals \((p)\). Primitive zeta orbits correspond to these
indecomposable elements.
\end{proof}

\begin{corollary}[Prime-orbit length]
The logarithmic length of a primitive finite zeta orbit is \(\log p\), where
\(p\) is the corresponding rational prime.
\end{corollary}

\begin{lemma}[Finite arithmetic trace]
\label{lem:FiniteArithmeticTrace}
For \(\Re(s)>1\), the finite prime-orbit trace is
\[
  \Tr_{\mathrm{fin}}(s)
  =
  \frac{\zeta'(s)}{\zeta(s)}.
\]
\end{lemma}

\begin{proof}
By \autoref{lem:PrimeOrbitCorrespondence}, the arithmetic trace-formula analogy
\cite{Connes1999} gives
\[
  \Tr_{\mathrm{fin}}(s)
  =
  -\sum_p\sum_{k\ge1}(\log p)e^{-ks\log p}
  =
  -\sum_p\sum_{k\ge1}(\log p)p^{-ks}.
\]
Equivalently,
\[
  \Tr_{\mathrm{fin}}(s)
  =
  -\sum_{n=1}^{\infty}\frac{\Lambda(n)}{n^s},
\]
where \(\Lambda\) is the von Mangoldt function
\cite{IwaniecKowalski2004}. The Euler product for \(\zeta\), valid for
\(\Re(s)>1\), gives the stated identity.
\end{proof}

\begin{lemma}[Completed zeta trace]
\label{lem:CompletedZetaTrace}
The logarithmic derivative of \(\mathcal E\) is
\[
  \frac{\mathcal E'(s)}{\mathcal E(s)}
  =
  -\frac12\log\pi
  +\frac12\psi\left(\frac{s}{2}\right)
  +\frac{\zeta'(s)}{\zeta(s)},
  \qquad
  \psi(z)=\frac{\Gamma'(z)}{\Gamma(z)}.
\]
\end{lemma}

\begin{proof}
Take the logarithmic derivative of
\[
  \mathcal E(s)
  =
  \frac12\,\pi^{-s/2}\Gamma\left(\frac{s}{2}\right)\zeta(s).
\]
\end{proof}

Putting \(s=\frac12+\sqrt{\lambda}\), \autoref{lem:CompletedZetaTrace} gives
\[
  \frac{\Phi'(\lambda)}{\Phi(\lambda)}
  =
  \frac{1}{2\sqrt{\lambda}}
  \left[
    -\frac12\log\pi
    +\frac12\psi\left(\frac14+\frac12\sqrt{\lambda}\right)
    +\frac{\zeta'(\frac12+\sqrt{\lambda})}
           {\zeta(\frac12+\sqrt{\lambda})}
  \right].
\]
For \(\Re(\frac12+\sqrt{\lambda})>1\), this becomes
\[
  \frac{\Phi'(\lambda)}{\Phi(\lambda)}
  =
  \frac{1}{2\sqrt{\lambda}}
  \left[
    -\frac12\log\pi
    +\frac12\psi\left(\frac14+\frac12\sqrt{\lambda}\right)
    -\sum_{n=1}^{\infty}
      \frac{\Lambda(n)}{n^{1/2+\sqrt{\lambda}}}
  \right].
\]

\begin{lemma}[Shifted-square equivalence]
\label{lem:ShiftedSquareEquivalence}
The Riemann hypothesis is equivalent to the assertion that every zero of
\(\Phi\) lies on \(\mathbb R_{\le0}\).
\end{lemma}

\begin{proof}
If the Riemann hypothesis holds, every zero of \(\mathcal E\) has the form
\[
  s=\frac12+i\gamma,
  \qquad \gamma\in\mathbb R.
\]
Thus
\[
  \lambda=\left(s-\frac12\right)^2=-\gamma^2\le0.
\]
Conversely, suppose every zero of \(\Phi\) lies on \(\mathbb R_{\le0}\). If
\(\rho\) is a zero of \(\mathcal E\), then
\[
  \lambda_\rho=\left(\rho-\frac12\right)^2=-\gamma^2
\]
for some \(\gamma\in\mathbb R\). Hence
\(\rho=\frac12\pm i\gamma\), so \(\Re(\rho)=\frac12\). Since the zeros of
\(\mathcal E\) are precisely the nontrivial zeros of \(\zeta(s)\), the Riemann
hypothesis follows \cite{Titchmarsh1986}.
\end{proof}

\begin{theorem}[Riemann hypothesis as tachyonic determinant zeros]
\label{thm:NoTachyonRH}
The Riemann hypothesis is equivalent to the statement that the determinant has
no tachyonic, or complex mass-squared, zeros: all determinant zeros occur at
\[
  \lambda=-m_n^2,
  \qquad
  m_n^2\in\mathbb R_{\ge0}.
\]
\end{theorem}

\begin{proof}
Under \autoref{def:ZetaDeterminant}, zeros of \(\Phi\) are precisely zeros of
the regularized determinant, hence occur at the negatives of the mass-squared
values encoded by \(M_\zeta^2\). By
\autoref{lem:ShiftedSquareEquivalence}, the Riemann hypothesis is equivalent to
all such zeros lying on \(\mathbb R_{\le0}\), which is exactly the absence of
tachyonic or complex mass-squared zeros in this determinant model.
\end{proof}

\begin{definition}[Arithmetic mass divisor]
Let \(\xi(s)\) denote the completed Riemann zeta function, and let
\[
Z_{\xi}:=\{\rho\in\mathbb C:\xi(\rho)=0,\ 0<\Re(\rho)<1\}
\]
be the nontrivial zero divisor, counted with multiplicity. Define the shifted
arithmetic mass-squared divisor
\[
\Sigma_{\zeta}
:=
\left\{
m_{\rho}^{2}:= -\left(\rho-\frac12\right)^{2}
\;:\;
\rho\in Z_{\xi}
\right\}.
\]
The functional equation identifies the pair \(\rho\) and \(1-\rho\), since
\[
-\left(1-\rho-\frac12\right)^2
=
-\left(\rho-\frac12\right)^2.
\]
Complex conjugation acts by
\[
m_{\rho}^{2}\longmapsto \overline{m_{\rho}^{2}}.
\]
We equip \(\Sigma_{\zeta}\) with its natural locally finite divisor topology.
\end{definition}

\begin{definition}[Arithmetic mass divisor]
Let \(\xi(s)\) denote the completed Riemann zeta function, and let
\[
Z_{\xi}:=\{\rho\in\mathbb C:\xi(\rho)=0,\ 0<\Re(\rho)<1\}
\]
be the nontrivial zero divisor, counted with multiplicity. Define the shifted
arithmetic mass-squared divisor
\[
\Sigma_{\zeta}
:=
\left\{
m_{\rho}^{2}:= -\left(\rho-\frac12\right)^{2}
\;:\;
\rho\in Z_{\xi}
\right\}.
\]
The functional equation identifies the pair \(\rho\) and \(1-\rho\), since
\[
-\left(1-\rho-\frac12\right)^2
=
-\left(\rho-\frac12\right)^2.
\]
Complex conjugation acts by
\[
m_{\rho}^{2}\longmapsto \overline{m_{\rho}^{2}}.
\]
We equip \(\Sigma_{\zeta}\) with its natural locally finite divisor topology.
\end{definition}

\begin{definition}[UCFT spectral algebra]
Let \(\phi_{0}\in \mathrm{EIII}\) be a vacuum and let
\[
\Delta_{\mathrm{UCFT}}
=
\nabla^{*}\nabla+\mathcal R+\operatorname{Hess}_{\phi_{0}}(V)
\]
be the quadratic fluctuation operator in the universal coset field theory.
After the required gauge fixing and removal of zero modes, take
\(\Delta_{\mathrm{UCFT}}\) to be a self-adjoint nonnegative operator on a
Hilbert space \(\mathcal H_{\mathrm{UCFT}}\). Let
\[
\Sigma_{\mathrm{UCFT}}
:=
\operatorname{Spec}(\Delta_{\mathrm{UCFT}})
\subseteq \mathbb R_{\ge0}.
\]
Define the commutative spectral \(C^{*}\)-algebra
\[
\mathcal A_{\mathrm{UCFT}}
:=
C_{0}(\Sigma_{\mathrm{UCFT}}).
\]
Similarly define
\[
\mathcal A_{\zeta}
:=
C_{0}(\Sigma_{\zeta}),
\]
where \(\Sigma_{\zeta}\) is the arithmetic mass-squared divisor.
\end{definition}

\begin{definition}[Galois-equivariant spectral realization]
A Galois-equivariant spectral realization of the zeta sector in the UCFT
sector is a \(*\)-isomorphism of commutative spectral \(C^{*}\)-algebras
\[
\Theta:\mathcal A_{\zeta}\xrightarrow{\sim}\mathcal A_{\mathrm{UCFT}}
\]
such that:
\begin{enumerate}
\item \(\Theta\) preserves the distinguished mass-squared coordinate;
\item \(\Theta\) intertwines complex conjugation on \(\Sigma_{\zeta}\) with the
real structure of the UCFT fluctuation Hilbert space;
\item at every finite algebraic approximation level, \(\Theta\) is compatible
with the relevant Galois action on the arithmetic mass divisor and the induced
algebraic symmetry action on the UCFT spectral data.
\end{enumerate}
\end{definition}

\begin{theorem}[Albert--EIII spectral criterion for the Riemann hypothesis]
For every Galois-equivariant spectral realization
\[
\Theta:\mathcal A_{\zeta}\xrightarrow{\sim}\mathcal A_{\mathrm{UCFT}},
\]
the Riemann hypothesis holds.
\end{theorem}

\begin{proof}
Since
\[
\Theta:\mathcal A_{\zeta}\xrightarrow{\sim}\mathcal A_{\mathrm{UCFT}}
\]
is a \(*\)-isomorphism of commutative \(C^{*}\)-algebras, Gelfand duality gives
a homeomorphism of spectra
\[
\Theta^{\vee}:\Sigma_{\mathrm{UCFT}}\xrightarrow{\sim}\Sigma_{\zeta}.
\]
The UCFT fluctuation operator is self-adjoint and nonnegative, hence
\[
\Sigma_{\mathrm{UCFT}}\subseteq\mathbb R_{\ge0}.
\]
Because the realization preserves the distinguished mass-squared coordinate,
the transported arithmetic mass-squared divisor also satisfies
\[
\Sigma_{\zeta}\subseteq\mathbb R_{\ge0}.
\]

Let \(\rho\in Z_{\xi}\). Write
\[
\rho=\frac12+\alpha+i\gamma,
\qquad
\alpha,\gamma\in\mathbb R.
\]
Then
\[
m_{\rho}^{2}
=
-\left(\rho-\frac12\right)^{2}
=
-(\alpha+i\gamma)^{2}
=
\gamma^{2}-\alpha^{2}-2i\alpha\gamma.
\]
Since \(m_{\rho}^{2}\in\mathbb R_{\ge0}\), its imaginary part vanishes:
\[
-2\alpha\gamma=0.
\]
For a nontrivial zero of the completed Riemann zeta function in the critical
strip, \(\gamma\neq0\). Hence
\[
\alpha=0.
\]
Thus
\[
\Re(\rho)=\frac12.
\]
Since this holds for every nontrivial zero \(\rho\), the Riemann hypothesis
follows.
\end{proof}

\begin{definition}[Locally finite divisor topology]
Let \(\operatorname{Div}_{\mathrm{loc}}(\mathbb C)\) denote the space of
locally finite effective divisors on \(\mathbb C\). Thus an element is a formal
sum
\[
D=\sum_{\lambda\in\mathbb C}m_{\lambda}[\lambda],
\qquad
m_{\lambda}\in\mathbb Z_{\ge0},
\]
such that only finitely many \(\lambda\)'s occur in each compact subset of
\(\mathbb C\).

We equip \(\operatorname{Div}_{\mathrm{loc}}(\mathbb C)\) with the vague
topology: a sequence \(D_n\) converges to \(D\) if and only if, for every
compactly supported continuous test function \(\varphi\in C_c(\mathbb C)\),
\[
\sum_{\lambda}m_{\lambda}^{(n)}\varphi(\lambda)
\longrightarrow
\sum_{\lambda}m_{\lambda}\varphi(\lambda).
\]
\end{definition}

\begin{definition}[Spectral continuity]
Let \(B\) be a parameter space containing a coset point \(b_{\mathrm{UCFT}}\)
and an arithmetic point \(b_{\zeta}\). A spectral bridge from the UCFT coset
sector to the zeta sector is a family of closed operators
\[
\{\mathcal M_b^2\}_{b\in B}
\]
such that:
\begin{enumerate}
\item each \(\mathcal M_b^2\) acts on a Hilbert space \(\mathcal H_b\);
\item the family is continuous in the norm-resolvent topology, meaning that
for some, equivalently every, \(w\in\mathbb C\setminus\mathbb R\),
\[
b\longmapsto (\mathcal M_b^2-w)^{-1}
\]
is norm-continuous;
\item each \(\mathcal M_b^2\) admits a regularized determinant in the
mass-squared coordinate,
\[
D_b(\mu):=\det_{\mathrm{reg}}(\mathcal M_b^2-\mu);
\]
\item the zero-divisor map
\[
b\longmapsto \operatorname{div}(D_b)
\]
is continuous as a map
\[
B\longrightarrow \operatorname{Div}_{\mathrm{loc}}(\mathbb C).
\]
\end{enumerate}
\end{definition}

\begin{definition}[Galois-compatible spectral bridge]
A spectral bridge is Galois-compatible if, at every finite algebraic
approximation level, the divisor
\[
\operatorname{div}(D_b)
\]
is stable under the relevant Galois action, and the transition maps in the
bridge intertwine this Galois action with the real structure on the quantum
Hilbert spaces \(\mathcal H_b\).

Equivalently, for every admissible finite approximation \(D_{b,N}\), every
Galois element \(\sigma\), and every point \(\lambda\) of the divisor, one has
\[
\sigma(\lambda)\in \operatorname{supp}(D_{b,N})
\]
with the same multiplicity, and this compatibility is preserved under passage
to the limiting divisor.
\end{definition}

\begin{definition}[Positive continuous Albert--zeta realization]
A positive continuous Albert--zeta realization is a Galois-compatible spectral
bridge
\[
\{\mathcal M_b^2\}_{b\in B}
\]
such that:
\begin{enumerate}
\item the UCFT endpoint satisfies
\[
\mathcal M_{b_{\mathrm{UCFT}}}^{2}
=
\Delta_{\mathrm{UCFT}};
\]
\item the arithmetic endpoint satisfies
\[
\operatorname{div}(D_{b_{\zeta}})
=
\Sigma_{\zeta}
=
\left\{
-\left(\rho-\frac12\right)^2:
\xi(\rho)=0,\ 0<\Re(\rho)<1
\right\};
\]
\item every operator in the bridge is self-adjoint and nonnegative:
\[
\mathcal M_b^2=\left(\mathcal M_b^2\right)^{*},
\qquad
\mathcal M_b^2\ge0
\]
for all \(b\in B\).
\end{enumerate}
\end{definition}

\begin{theorem}[Continuous spectral criterion for the Riemann hypothesis]
Every positive continuous Albert--zeta realization implies the Riemann
hypothesis.
\end{theorem}

\begin{proof}
For each \(b\in B\), self-adjointness implies
\[
\operatorname{Spec}(\mathcal M_b^2)\subseteq\mathbb R.
\]
Nonnegativity strengthens this to
\[
\operatorname{Spec}(\mathcal M_b^2)\subseteq\mathbb R_{\ge0}.
\]
Since
\[
D_b(\mu)=\det_{\mathrm{reg}}(\mathcal M_b^2-\mu),
\]
the zero divisor of \(D_b\) lies in the mass-squared spectrum:
\[
\operatorname{div}(D_b)\subseteq\mathbb R_{\ge0}.
\]

At the arithmetic endpoint \(b=b_{\zeta}\), the zero divisor is
\[
\Sigma_{\zeta}
=
\left\{
-\left(\rho-\frac12\right)^2:
\xi(\rho)=0,\ 0<\Re(\rho)<1
\right\}.
\]
Thus every shifted arithmetic mass-squared value is real and nonnegative.

Let
\[
\rho=\frac12+\alpha+i\gamma,
\qquad
\alpha,\gamma\in\mathbb R.
\]
Then
\[
-\left(\rho-\frac12\right)^2
=
-(\alpha+i\gamma)^2
=
\gamma^2-\alpha^2-2i\alpha\gamma.
\]
Since this number is real, its imaginary part vanishes:
\[
-2\alpha\gamma=0.
\]
A nontrivial zero of \(\xi(s)\) in the critical strip is not real, so
\[
\gamma\neq0.
\]
Hence
\[
\alpha=0.
\]
Therefore
\[
\Re(\rho)=\frac12.
\]
Since this holds for every nontrivial zero, the Riemann hypothesis follows.
\end{proof}

%======================================================================%
\section{Heterotic Interface and Local Matching}
\label{sec:Holographic}
%======================================================================%

The construction developed above realizes gravity and spacetime as
\emph{induced} structures: there is no fundamental graviton, and the metric
$g_{\mu\nu}$ is the pull-back of the EIII metric along a spacetime-filling
configuration (\autoref{thm:EmergentRiemannian}), while the Planck mass is
generated radiatively in the sense of Sakharov and Adler
\cite{SakharovInducedGravity,AdlerInducedGravity}. This motivates a precise,
but limited, comparison between the coset sector and the massless bosonic
sector of an $E_{8}\times E_{8}$ heterotic compactification \cite{GHMR:1985}
whose visible gauge group is $\ESix$. We emphasize at the outset that the
statement established here is local, two-derivative, and current-algebraic; it
is an interface with a heterotic visible sector, not a proof of an exact
holographic duality.

\subsection{The heterotic datum}

Fix an $E_{8}\times E_{8}$ heterotic background on
$\mathbb R\smexp{1,3}\times X$ with $X$ a six-dimensional internal space whose
worldsheet realization preserves an $\ESix$ current algebra at level one, so
that the visible massless spectrum organizes into $\ESix$ multiplets and the
$78$ adjoint is present as a gauge multiplet.  Let $M_{\mathrm{res}}$ denote the
residual scalar target obtained after the internal CFT has fixed all moduli
except those of the visible $\ESix$-breaking sector.  We write $H=\SOTenXUOne_\psi$
for the grand-unified isotropy group singled out by an adjoint Higgs vacuum in
the $\ESix\supset H$ singlet direction (\autoref{sec:Z12}).

\begin{definition}[Interface map]
\label{def:Interface}
The \emph{interface map} is the $\ESix$-equivariant map
\[
  F\ :\ M_{\mathrm{res}}\;\longrightarrow\;\M=\MCosetInline
\]
sending the gauge-unfixed adjoint-Higgs orbit of the heterotic model onto
$\M$, and identifying the residual scalar fluctuations with the
Goldstone coordinates $\phi\smexp{\alpha}$ of \autoref{sec:EmergentFields}.
\end{definition}

\begin{lemma}[Two-derivative matching]
\label{lem:TwoDerivative}
On the image of $F$ the heterotic non-linear sigma-model metric restricted to the
visible $\ESix$-breaking directions coincides, up to an overall normalization
$f\smexp{2}$, with the pull-back $F\smexp{\ast}\mathfrak g$ of the unique
$\ESix$-invariant EIII metric $\mathfrak g_{\alpha\beta}=f\smexp{2}(-K_{\alpha\beta})$.
\end{lemma}

\begin{proof}
The residual visible directions transform in $16_{-3}\oplus\overline{16}_{+3}$
under $H$, which is real-irreducible (\autoref{thm:EmergentRiemannian}).  Any
$H$-invariant symmetric two-tensor on this representation is unique up to scale
by Schur's lemma; both the sigma-model metric and $\mathfrak g$ are such
tensors, hence proportional.  The proportionality constant defines
$f\smexp{2}$.
\end{proof}

\begin{lemma}[Current-algebra matching]
\label{lem:Current}
The $H=\SOTenXUOne_\psi$ isometry action on $\M$ is generated by the same
currents as the level-one $\SOTen$ and $\UOne_\psi$ Ka\v{c}--Moody currents of
the heterotic worldsheet theory, with matched central data.  In particular the
$\UOne_\psi$ charge assignment $16_{-3}\oplus\overline{16}_{+3}$ reproduces the
worldsheet charge lattice of the visible matter.
\end{lemma}

\begin{proof}
The Killing vector fields of $\M$ generating $H$ act on the tangent
representation exactly as the zero modes of the $\SOTen\times\UOne_\psi$
currents act on the visible states; the charges agree by construction of $F$,
and the $\UOne_\psi$ normalization is fixed by requiring the $16$ to carry
charge $-3$, matching the branching $78=45_{0}\oplus1_{0}\oplus16_{-3}\oplus\overline{16}_{+3}$.
\end{proof}

\begin{theorem}[Local heterotic realization]
\label{thm:Holographic}
Let $F$ be the interface map of \autoref{def:Interface}.  Then, in a
neighbourhood of the gauge-unfixed Higgs orbit, the two-derivative bosonic
action of the visible $\ESix$-breaking sector of the heterotic model is mapped by
$F$ onto the coset action of \autoref{sec:EmergentFields}, with the induced
Einstein--Hilbert term of \autoref{sec:sigma} providing the gravitational
sector.  Under this map the Green--Schwarz axion \cite{GreenSchwarz:1984} of the
heterotic background is identified with the $\UOne_\psi$ Goldstone mode of $\M$.
\end{theorem}

\begin{proof}
Combine \autoref{lem:TwoDerivative} and \autoref{lem:Current}: the kinetic term
matches by two-derivative matching and the symmetry action matches by
current-algebra matching, so the visible bosonic Lagrangians agree to two
derivatives.  The gravitational sector is not present in the heterotic
two-derivative visible action as an independent field; on the side it is
induced (\autoref{sec:sigma}), consistent with the absence of a fundamental
visible-sector graviton.  The Green--Schwarz axion couples to
$\Tr F\wedge F$ with a shift symmetry, and the interface identifies the
corresponding visible shift mode with the $\UOne_\psi$ Goldstone of $\M$.
\end{proof}

\begin{table}[t]
\centering
\caption{Interface dictionary between the visible heterotic sector and the
coset sector. The correspondence is established at the level of the
two-derivative bosonic action and the $\SOTen\times\UOne_\psi$ current algebra.}
\label{tab:Dictionary}
\begin{tabular}{@{}ll@{}}
\toprule
Heterotic $E_{8}\times E_{8}$ visible sector & coset sector \\
\midrule
$\ESix$ current algebra (level one)          & Compact $\ESix$ symmetry selected by $\jalg$ \\
Adjoint Higgs orbit $E_{6}\cdot\langle78_H\rangle$ & Vacuum manifold $\M=\MCosetInline$ \\
Visible matter $16_{-3}\oplus\overline{16}_{+3}$    & Goldstone tangent space $T\M$ \\
Sigma-model metric on visible moduli          & EIII metric $\mathfrak g=f\smexp{2}(-K)$ \\
$\SOTen\times\UOne_\psi$ Ka\v{c}--Moody currents & Isometry generators of $\M$ \\
Green--Schwarz axion                          & $\UOne_\psi$ Goldstone mode \\
No fundamental visible graviton               & Induced (Sakharov--Adler) gravity \\
\bottomrule
\end{tabular}
\end{table}

\begin{remark}[Scope of the correspondence]
\label{rmk:HolographicScope}
\autoref{thm:Holographic} is a statement about the two-derivative bosonic action
and the current algebra only. It does not assert an exact conformal duality, a
matching of the full string spectrum, or agreement of higher-derivative
$\alpha'$ corrections. In particular the heterotic side carries an infinite
tower of massive string states with no coset counterpart, and the coset side is an
effective field theory valid below the string scale. The value of the
correspondence is structural: it exhibits the coset as the low-energy image
of a genuine string background, and it shows that the induced-gravity mechanism
is compatible with a formulation in which no fundamental visible-sector
graviton is postulated.
\end{remark}


%======================================================================%
\section{The \texorpdfstring{$\mathbb Z_{12}$}{Z12} Asymmetric-Orbifold Realization}
\label{sec:Z12}
%======================================================================%

We now consider a heterotic compactification intended to meet the criteria of
\autoref{sec:Holographic}: a $\mathbb Z_{12}$ asymmetric orbifold of the
$E_{8}\times E_{8}$ heterotic string \cite{DHVW:1985,Narain:1986} whose visible
sector is a three-generation $\ESix$ grand unified theory with an adjoint Higgs
$78_{H}$ \cite{Z12Heterotic}.  Giving the adjoint Higgs a vacuum expectation
value in the $\ESix\supset H$ singlet direction breaks $\ESix\to H=\SOTenXUOne_\psi$,
and the local orbit computation identifies the resulting gauge-unfixed Higgs
orbit with the coset $\M$.

\subsection{The orbifold datum}

The model is a $\mathbb Z_{12}$ asymmetric orbifold with visible gauge group
$\ESix$ realized at level one, three chiral generations in the $27$, and an
adjoint Higgs multiplet $78_{H}$ arising from the diagonal embedding of the
visible $\ESix$ in the ten-dimensional gauge group.  Under
$\ESix\supset H=\SOTenXUOne_\psi$ the adjoint branches as
\begin{equation}
  78 \;=\; 45_{0}\,\oplus\,1_{0}\,\oplus\,16_{-3}\,\oplus\,\overline{16}_{+3},
  \label{eq:AdjointBranch}
\end{equation}
and we take the Higgs vacuum $\langle78_{H}\rangle=\Phi_{0}$ to lie in the
$1_{0}$ direction.

\begin{theorem}[Gauge-unfixed EIII orbit theorem]
\label{thm:Z12Orbit}
Let $\Phi_{0}$ be an adjoint-Higgs vacuum in the $1_{0}\subset78$ direction under
$\ESix\supset H=\SOTenXUOne_\psi$.  Then the gauge-unfixed Higgs orbit is
\[
  \mathcal O_{H}
  \;=\;\ESix\cdot\Phi_{0}
  \;\cong\;\MCosetInline
  \;=\;\M,
  \qquad
  T_{\mathbb C}\,\mathcal O_{H}=16_{-3}\oplus\overline{16}_{+3}.
\]
\end{theorem}

\begin{proof}
By the branching \eqref{eq:AdjointBranch} the centralizer in $\ESix$ of a
generic vector in the $1_{0}$ direction is exactly $H=\SOTenXUOne_\psi$: the
$45_{0}$ commutes with $\Phi_{0}$, the $\UOne_\psi$ generator is $\Phi_{0}$
itself, and the $16_{-3}\oplus\overline{16}_{+3}$ carry non-zero $\UOne_\psi$
charge and are therefore rotated.  By the orbit--stabilizer theorem
$\ESix\cdot\Phi_{0}\cong\ESix/H=\MCosetInline$.  The tangent space at $\Phi_{0}$
is the span of the broken generators, namely
$16_{-3}\oplus\overline{16}_{+3}$, giving the stated complexified tangent
representation.  This is the coset $\M$ of \autoref{sec:EmergentFields}.
\end{proof}

\begin{theorem}[Gauge-fixed/gauge-unfixed equivalence]
\label{thm:Z12Equivalence}
The EIII scalar sector is physically equivalent to the gauge-unfixed
Stueckelberg/Higgs-orbit description of the $\ESix\to\SOTenXUOne_\psi$ breaking
sector of the $\mathbb Z_{12}$ model.  Unitary gauge maps this description to the
standard massive-vector description of the heterotic $\ESix$ grand unified
theory, in which the $16_{-3}\oplus\overline{16}_{+3}$ Goldstones are absorbed by
the broken gauge bosons.
\end{theorem}

\begin{proof}
In the gauge-unfixed description the $32$ real coordinates of
$\mathcal O_{H}\cong\M$ are Stueckelberg/Goldstone fields for the broken
$\ESix/H$ directions.  Fixing unitary gauge sets the Goldstones to zero and gives
mass to the corresponding gauge bosons via the Higgs mechanism; the massive
vector content is gauge-invariant and hence the same in both descriptions. The
coset action of \autoref{sec:EmergentFields} is the gauge-unfixed
Goldstone dynamics on $\mathcal O_{H}$, so the two descriptions are equivalent
at the level of this sector.
\end{proof}

\begin{corollary}[Breaking chain]
\label{cor:BreakingChain}
The $\mathbb Z_{12}$ model realizes the phenomenological breaking chain
\[
  \ESix\;\to\;\SOTenXUOne_\psi\;\to\;\SU5\;\to\;\SMSubgroup,
\]
with the first step realized geometrically as the emergence of the coset
$\M$ and the subsequent steps by the standard heterotic Higgs mechanism.
\end{corollary}

\begin{proof}
The first arrow is \autoref{thm:Z12Orbit}.  The remaining arrows are the
standard $\SOTen\to\SU5\to\SMSubgroup$ descent of a heterotic $\ESix$ grand
unified theory with the matter content of the $27$; three generations follow
from the three fixed-point families of the $\mathbb Z_{12}$ action
\cite{Z12Heterotic}.
\end{proof}

\begin{corollary}[UV completion]
\label{cor:UVComplete}
If the $\mathbb Z_{12}$ model of \cite{Z12Heterotic} satisfies the global
heterotic consistency conditions listed in \autoref{rmk:Ledger}, then it is an
ultraviolet completion of the effective theory in the sense of the local
correspondence of \autoref{sec:Holographic}.
\end{corollary}

\begin{proof}
Immediate from \autoref{thm:Z12Orbit}, \autoref{thm:Holographic}, and the
definition of the global heterotic consistency conditions: local matching plus
global worldsheet consistency is the required condition for the heterotic model
to define a consistent string background whose low-energy visible sector is the
coset theory.
\end{proof}


%======================================================================%
\section{The Albert--EIII Swampland Criterion}
\label{sec:Swampland}
%======================================================================%

The results above give a \emph{relative} swampland criterion. We do not claim
that all non-coset vacua are inconsistent. The statement is restricted to the
Albert-selected heterotic class: within that class, a vacuum can realize the
coset correspondence only when its gauge-unfixed Higgs/moduli sector is EIII with
the data specified below.

\begin{theorem}[Albert--EIII swampland theorem]
\label{thm:Swampland}
Let $\mathcal T$ be a four-dimensional effective theory proposed as the
Albert-selected low-energy sector of an $E_{8}\times E_{8}$ heterotic
compactification.  Assume that its visible exceptional sector is selected by the
Albert algebra $\jalg$ and that the corresponding compact real form is $\ESix$.
Then $\mathcal T$ lies in the Albert-selected heterotic landscape only if its
gauge-unfixed Higgs/moduli sector contains a residual scalar orbit
\[
  \mathcal O\;\cong\;\MCosetInline\;=\;\M,
  \qquad
  T_{\mathbb C}\,\mathcal O=16_{-3}\oplus\overline{16}_{+3},
\]
carrying the unique positive $\ESix$-invariant EIII metric, globally consistent
$\SOTenXUOne_\psi$ gauge data, anomaly cancellation, spin-lift compatibility,
Green--Schwarz matching, and heterotic worldsheet consistency.  Equivalently, any
proposed Albert-selected four-dimensional effective theory failing one of these
conditions lies outside the Albert--EIII landscape: it has no
$E_{8}\times E_{8}$ heterotic ultraviolet completion of the specified type
realizing the coset correspondence.
\end{theorem}

\begin{proof}
The argument is by necessity together with the contrapositive.

\emph{Type and real form.}  Albert selection fixes the complex exceptional type
to $E_{6}$ (\autoref{sec:AxiomsAlbertSSB}); choosing the compact real form gives the
compact $\ESix$ structure used in the EIII target and guarantees a
positive-definite invariant metric (\autoref{thm:uniqueMetric}).

\emph{Forced orbit.}  Imposing the grand-unified isotropy
$H=\SOTenXUOne_\psi$ forces the homogeneous scalar orbit to be $\ESix/H$.  By
the orbit--stabilizer theorem an adjoint Higgs vacuum in the $1_{0}$ direction of
the branching \eqref{eq:AdjointBranch} has centralizer $H$, so its gauge-unfixed
Higgs orbit is exactly $\ESix\cdot\Phi_{0}\cong\MCosetInline$
(\autoref{thm:Z12Orbit}).

\emph{Rigidity.}  The tangent representation is forced to be
$16_{-3}\oplus\overline{16}_{+3}$, which is real-irreducible.  Hence the
$H$-invariant positive metric is unique up to scale, and there is no competing
compact homogeneous residual scalar sector with the same Albert/$\ESix$/$H$
data (\autoref{thm:EmergentRiemannian}).

\emph{No-soldering exclusion.} Because $\mathfrak m$ has no invariant
four-plane and the EIII metric is positive definite, the coset cannot supply a
Lorentzian vierbein by selecting four invariant scalar directions. The valid
construction is the nondegenerate Euclidean pull-back of
\autoref{thm:EmergentRiemannian}, followed by Osterwalder--Schrader
reconstruction; gravity then enters through the induced Sakharov--Adler
dynamics of \autoref{sec:sigma}.

\emph{Global constraints.}  Local EIII matching is not sufficient; a heterotic
completion must further satisfy modular invariance, GSO/fermion-parity
consistency, flux and Wilson-line quantization, bundle transition-function
compatibility, anomaly cancellation on overlaps, spin-lift conditions for the
$16$ matter, and Green--Schwarz matching \cite{GreenSchwarz:1984}.

\emph{Contrapositive.}  If a proposed effective theory lacks the EIII orbit, the
required tangent representation, the EIII metric, the correct gauge action, the
anomaly/Green--Schwarz compatibility, the spin-lift data, or global heterotic
consistency, then it cannot be the Albert-selected $E_{8}\times E_{8}$
ultraviolet completion of the specified type.
\end{proof}

\begin{remark}[Global heterotic consistency ledger]
\label{rmk:Ledger}
The global heterotic consistency conditions invoked above and below are:
modular invariance, GSO/fermion-parity consistency, flux and Wilson-line
quantization, bundle transition-function compatibility, anomaly cancellation
on overlaps, spin-lift conditions for the $16$ matter, and Green--Schwarz
matching \cite{GreenSchwarz:1984}.
\end{remark}

\begin{corollary}[$\mathbb Z_{12}$ test]
\label{cor:Z12Test}
A $\mathbb Z_{12}$ asymmetric heterotic $\ESix$ orbifold realizes the coset
correspondence if its gauge-unfixed adjoint-Higgs orbit is $\MCosetInline$ and
if its full worldsheet, anomaly, bundle, flux, spectrum, renormalization-group,
and gravitational matching data satisfy the global heterotic consistency
conditions of \autoref{rmk:Ledger}.  If these checks fail, the model remains an
$\ESix$-aligned heterotic model but not the coset ultraviolet completion.
\end{corollary}

\begin{proof}
The orbit condition is \autoref{thm:Z12Orbit} and places the model in the
Albert-selected landscape of \autoref{thm:Swampland}; the global conditions are
those of \autoref{rmk:Ledger} and \autoref{cor:UVComplete}.  If any global
condition fails, the contrapositive of \autoref{thm:Swampland} applies.
\end{proof}

\begin{corollary}[$\mathbb Z_{12}$ candidate corollary]
\label{cor:Z12Leading}
The $\mathbb Z_{12}$ asymmetric $\ESix$ orbifold is a concrete candidate for
realizing the coset correspondence because it supplies an $E_{8}\times E_{8}$
heterotic origin, a visible $\ESix$ grand unified sector, three generations,
and an adjoint Higgs obtained by diagonal embedding. If its gauge-unfixed
adjoint-Higgs orbit, anomaly data, worldsheet data, and gravitational reduction
satisfy the global heterotic equivalence criterion of \autoref{rmk:Ledger},
then it gives a heterotic ultraviolet completion of the coset theory.
\end{corollary}

\begin{proof}
The four structural features are those of the model of
\cite{Z12Heterotic} recorded in \autoref{sec:Z12}; each is a prerequisite of
\autoref{thm:Swampland}, and no other presently identified construction is known
to supply all four together with the EIII orbit of \autoref{thm:Z12Orbit}.
\end{proof}

\begin{remark}[Status of the criterion]
\label{rmk:SwamplandStatus}
\autoref{thm:Swampland} is a relative, class-restricted statement: it constrains
only theories proposed as Albert-selected heterotic vacua. It is
not a universal assertion that all other four-dimensional effective theories are
inconsistent. Within its stated scope it is sharp, and it makes the
$\mathbb Z_{12}$ model of \autoref{sec:Z12} a concrete landscape candidate
subject to the checklist of \autoref{cor:Z12Test}.
\end{remark}


%======================================================================%
\section*{Declarations}
%======================================================================%

The author declares no competing interests. No datasets were generated or
analyzed in this work. No external funding is reported.

\begin{appendices}


\appendix

%======================================================================%
\section{Albert Algebra Proofs}
\label{app:Albert}
%======================================================================%

This appendix supplies proofs regarding the Albert algebra deferred from \autoref{sec:AxiomsAlbertSSB}.

%%%%%%%%%%%%%%%%%%%%%%%%%%%%%%%%%%%%%%%%%%%%%%%%%%%%%%%%%%%%%%%%%%%%%%
\subsection{Dimension and Jordan Identity}
\label{app:Albert:JordanDim}
%%%%%%%%%%%%%%%%%%%%%%%%%%%%%%%%%%%%%%%%%%%%%%%%%%%%%%%%%%%%%%%%%%%%%%

\begin{proof}[Proof of \autoref{thm:JordanDim}]
Choose a basis $\{e_{0},e_{1},\dots,e_{7}\}$ of the octonions satisfying $e_{0}=1$ and $e_{i}\smexp{2}=-1$ for $i\ge1$.  
Every $X\in\jalg$ can be written uniquely as the Hermitian $3\times3$ matrix
\[
  X=\begin{pmatrix}
        \alpha_{1} & x_{3} & \bar x_{2}\\
        \bar x_{3} & \alpha_{2} & x_{1}\\
        x_{2} & \bar x_{1} & \alpha_{3}
      \end{pmatrix},
  \quad
  \alpha_{i}\in\mathbb R,
  \;
  x_{i}\in\Oct.
\]
The real-dimension count is therefore
$3\text{ scalars}+3\times8\text{ octonion components}=27$.
Jordan commutativity is obvious from the symmetric product $\circ$.  
The Jordan identity
$a\circ(b\circ a\smexp{2})=(a\circ b)\circ a\smexp{2}$
holds for every special Jordan algebra; it follows from the alternative law of the octonions \cite{McCrimmon:2004}.  Hence $\jalg$ is indeed a $27$-dimensional Jordan algebra.
\end{proof}

%%%%%%%%%%%%%%%%%%%%%%%%%%%%%%%%%%%%%%%%%%%%%%%%%%%%%%%%%%%%%%%%%%%%%%
\subsection{Quadratic and Cubic Invariants}
\label{app:Albert:Invs}
%%%%%%%%%%%%%%%%%%%%%%%%%%%%%%%%%%%%%%%%%%%%%%%%%%%%%%%%%%%%%%%%%%%%%%

\begin{proof}[Proof of \autoref{thm:Invs}]
For $X$ as above set
\[
  Q(X)=\Tr(X\smexp{2}), 
  \quad
  N(X)=\alpha_{1}\alpha_{2}\alpha_{3}-\sum_{i=1}\smexp{3}\alpha_{i}\|x_{i}\|\smexp{2}
        +2\,\Re\bigl(x_{1}x_{2}x_{3}\bigr),
\]
where juxtaposition is octonion multiplication. The form $Q$ is quadratic and $N$ is cubic in the real coordinates. Direct calculation shows
$
  Q(g \cdot X)=Q(X)
$
and
$
  N(g \cdot X)=\lambda(g)N(X)
$
for every $g\in\Str(\jalg)$ with multiplicative character $\lambda(g)\in\mathbb R\smexp{\times}$ \cite{Jacobson:1968}. In particular, the subgroup with $\lambda(g)=1$ preserves both forms.
\end{proof}

%%%%%%%%%%%%%%%%%%%%%%%%%%%%%%%%%%%%%%%%%%%%%%%%%%%%%%%%%%%%%%%%%%%%%%
\subsection{Exceptional Type \texorpdfstring{$E_{6}$}{E6}}
\label{app:Albert:E6}
%%%%%%%%%%%%%%%%%%%%%%%%%%%%%%%%%%%%%%%%%%%%%%%%%%%%%%%%%%%%%%%%%%%%%%

\begin{proof}[Proof of \autoref{thm:E6}]
Let $\Str\smexp{0}(\jalg)$ denote the connected component of the identity in
$\Str(\jalg)$. Jacobson's construction \cite{Jacobson:1968} realizes the
exceptional structure algebra as the sum of derivations of $\jalg$ (an
$\mathfrak f_{4}$) plus the traceless right multiplications
\[
  R_{X}: Y\mapsto X\circ Y-\frac13 Q(X,Y)\mathbf 1.
\]
After complexification this Lie algebra has type $\mathfrak e_{6}$. The main
text uses the simply connected compact real form of the same complex Lie type
for the EIII homogeneous space; this is the group denoted $\ESix$ throughout
the coset construction.
\end{proof}

%%%%%%%%%%%%%%%%%%%%%%%%%%%%%%%%%%%%%%%%%%%%%%%%%%%%%%%%%%%%%%%%%%%%%%
\subsection{Orbit Classification}
\label{app:Albert:OrbitClass}
%%%%%%%%%%%%%%%%%%%%%%%%%%%%%%%%%%%%%%%%%%%%%%%%%%%%%%%%%%%%%%%%%%%%%%

\begin{lemma}\label{lem:OrbitInv}
If $g\in \ESix$ and $X\in\jalg$ then
$Q(g\cdot X)=Q(X)$ and $N(g\cdot X)=N(X)$.
\end{lemma}

\begin{proof}
The compact $\ESix$ action used on the regular stratum preserves the restricted
quadratic and cubic invariants, so the claim follows from \autoref{thm:Invs}.
\end{proof}

\begin{proof}[Proof of \autoref{thm:OrbitClass}]
Suppose $g \cdot X=Y$. \autoref{lem:OrbitInv} gives equality of $(Q,N)$.
Conversely, restrict to the regular compact stratum specified in the theorem.
The standard invariant-theoretic description of the compact $E_{6}$ action
\cite{Springer:2000} shows that the common level set of $(Q,N)$ through
$X_{0}$ is a single compact orbit. Hence two points on that level set differ by
an element of $\ESix$.
\end{proof}

%%%%%%%%%%%%%%%%%%%%%%%%%%%%%%%%%%%%%%%%%%%%%%%%%%%%%%%%%%%%%%%%%%%%%%
\subsection{An \texorpdfstring{$\SOTenXUOne$}{SO(10) x U(1)} Stabilizer}
\label{app:Albert:Isotropy}
%%%%%%%%%%%%%%%%%%%%%%%%%%%%%%%%%%%%%%%%%%%%%%%%%%%%%%%%%%%%%%%%%%%%%%

\begin{proof}[Proof of \autoref{cor:Isotropy}]
The compact group of type $E_{6}$ has the maximal-rank subgroup
$\Spin(10)\times\UOne$ \cite{Slansky:1981}. Choosing $X_{0}$ in the
corresponding regular orbit gives stabilizer Lie algebra
$\mathfrak{so}(10)\oplus\mathfrak u(1)$ of dimension $45+1=46$. With the
global convention used in this paper, the stabilizer is denoted
$\SOTenXUOne$.
\end{proof}

%%%%%%%%%%%%%%%%%%%%%%%%%%%%%%%%%%%%%%%%%%%%%%%%%%%%%%%%%%%%%%%%%%%%%%
\subsection{Coercivity of the Symmetry-Breaking Potential}
\label{app:Albert:Coercive}
%%%%%%%%%%%%%%%%%%%%%%%%%%%%%%%%%%%%%%%%%%%%%%%%%%%%%%%%%%%%%%%%%%%%%%

\begin{lemma}\label{lem:A-coercive}
For $\alpha>0$, $k\ge2$ even, the polynomial
$V(X)=P(X)+\alpha[I_{2}(X)-c_{2}]\smexp{k}$ satisfies
$V(X)\to\infty$ as $\|X\|\to\infty$.
\end{lemma}

\begin{proof}
Write $I_{2}(X)=Q(X)=\beta\|X\|\smexp{2}+O(\|X\|)$ for some $\beta>0$.
Since $k$ is even,
\[
  \alpha[I_{2}(X)-c_{2}]\smexp{k}
  =\alpha\beta\smexp{k}\|X\|\smexp{2k}
   +O\bigl(\|X\|\smexp{2k-1}\bigr),
\]
so $V$ is radially unbounded.
\end{proof}

\begin{lemma}\label{lem:A-sumsq}
$P(X)\ge0$ for all $X\in\jalg$, with equality if and only if
$X\in \ESix \cdot X_{0}$ on the regular stratum.
\end{lemma}

\begin{proof}
Non-negativity is clear from \eqref{eq:Pdef}.
Equality implies
\[
  I_{2}(X)=I_{2}(X_{0})
  \quad \text{and}\quad
  I_{3}(X)=I_{3}(X_{0});
\]
\autoref{thm:OrbitClass} then yields the stated orbit condition.
\end{proof}

%======================================================================%
\section{Emergent Spacetime Proofs}
\label{app:Lorentz}
%======================================================================%

This appendix supplies proofs for \autoref{lem:Rank4} and \autoref{thm:EmergentRiemannian} from \autoref{sec:EmergentFields}.

%%%%%%%%%%%%%%%%%%%%%%%%%%%%%%%%%%%%%%%%%%%%%%%%%%%%%%%%%%%%%%%%%%%%%%
\subsection{Rank-Four Embedding}
\label{app:Lorentz:Rank4}
%%%%%%%%%%%%%%%%%%%%%%%%%%%%%%%%%%%%%%%%%%%%%%%%%%%%%%%%%%%%%%%%%%%%%%

Recall that the broken directions of the coset
\[
    \M = \MCoset
\]
form a $32$-dimensional real vector space
\(\mathfrak m\) carrying the unique invariant metric
\(\mathfrak g_{\alpha\beta}=f\smexp{2}(-K_{\alpha\beta})\) of
\autoref{thm:uniqueMetric}, which on the compact real form of \(\ESix\) is
positive definite. Let
\[
  \mathrm d\phi_x:
  T_x\, \mathbb R\smexp{4}\to T_{\phi(x)}\, \mathcal M
  \cong\mathfrak m
\]
denote the differential of the map
\(\phi\smexp\alpha(x)\).

\begin{proof}[Proof of \autoref{lem:Rank4}]
The pull-back at \(x\) is the quadratic form
\[
  g_{\mu\nu}(x)
  =\mathfrak g_{\alpha\beta}\,\partial_\mu\phi\smexp{\alpha}\partial_\nu\phi\smexp{\beta}
  =\mathfrak g\bigl(\mathrm d\sm\phi_x(\partial_\mu),\,
          \mathrm d\sm\phi_x(\partial_\nu)\bigr).
\]
Because \(\mathfrak g\) is positive definite, the associated Gram map obeys
\(\operatorname{rank} g_{\mu\nu}=\operatorname{rank}\mathrm d\phi_x=\dim
V_x\), where \(V_x=\Im\,\mathrm d\phi_x\subseteq\mathfrak m\): indeed
\(g_{\mu\nu}v\smexp{\mu}=0\) iff \(\mathfrak g(\mathrm d\phi_x v,\,\cdot\,)=0\)
iff \(\mathrm d\phi_x v=0\) by definiteness. Since
\(\dim T_x\mathbb R\smexp{4}=4\) we have \(\dim V_x\le4\), and non-degeneracy of
\(g_{\mu\nu}\) forces \(\dim V_x=4\), i.e.\ \(\mathrm d\phi_x\) is injective and
\(V_x=\mathfrak p_{0}(x)\) is a four-plane. Its \(\mathfrak g\)-orthogonal
complement \(\mathfrak p_{m}(x)\) has dimension \(28\).
\end{proof}

%%%%%%%%%%%%%%%%%%%%%%%%%%%%%%%%%%%%%%%%%%%%%%%%%%%%%%%%%%%%%%%%%%%%%%
\subsection{Euclidean Signature and Continuation}\label{app:Lorentz:Lorentz}
%%%%%%%%%%%%%%%%%%%%%%%%%%%%%%%%%%%%%%%%%%%%%%%%%%%%%%%%%%%%%%%%%%%%%%

\begin{proof}[Proof of \autoref{thm:EmergentRiemannian}]
Embed \(\SOTenXUOne\) in \(\ESix\) via the stabilizer of
\(X_{0}=\operatorname{diag}(1,1,1)\in\jalg\), and write
\(\mathfrak e_{6}=\mathfrak h\oplus\mathfrak m\) with
\(\mathfrak h=\mathfrak{so}(10)\oplus\mathfrak u(1)\). On the \emph{compact}
real form of \(\ESix\) the Killing form is negative definite on all of
\(\mathfrak e_{6}\); hence \(-K\), and therefore
\(\mathfrak g=f\smexp{2}(-K)\), is positive definite on \(\mathfrak m\). (This
is the same fact that makes the symmetric space
\(\M=\MCoset\) Riemannian; EIII is Hermitian symmetric, hence Riemannian, and
carries no invariant indefinite metric.)

By \autoref{lem:Rank4} the tetrad
\(e_\mu\smexp{a}=\theta\smexp{a}_{\alpha}\partial_\mu\phi\smexp{\alpha}\) is
invertible for a non-degenerate configuration. Choosing a
\(\mathfrak g\)-orthonormal basis \(\{u_a\}_{a=0}\smexp{3}\) of
\(\mathfrak p_{0}(x)\), so that \(\mathfrak g(u_a,u_b)=\delta_{ab}\), and
decomposing \(\partial_\mu\phi\smexp{\alpha}=e_\mu\smexp{a}u_a\smexp{\alpha}\),
\[
  g_{\mu\nu}
  =\mathfrak g(u_a,u_b)\,e_\mu\smexp{a}e_\nu\smexp{b}
  =\delta_{ab}\,e_\mu\smexp{a}e_\nu\smexp{b},
\]
which is symmetric, non-degenerate (invertibility of \(e_\mu\smexp{a}\)), and
positive definite. Smoothness in \(x\) follows from smoothness of \(\phi\) and
of the Maurer-Cartan form. Openness and density of the non-degenerate locus
hold because \(\det e_\mu\smexp{a}\neq0\) is an open condition satisfied by
generic immersions. Thus \(g_{\mu\nu}\) is a Riemannian metric on
\(\mathbb R\smexp{4}\), namely the first fundamental form of \(\phi\).

The signature is \((4,0)\), not \((1,3)\): a pull-back of a positive-definite
form is positive semidefinite for every \(\phi\), so no choice of configuration
can produce a Lorentzian induced metric.
\end{proof}

%======================================================================%
\section{Heat-Kernel Coefficients and Induced Gravity}
\label{app:Gravity}
%======================================================================%

%%%%%%%%%%%%%%%%%%%%%%%%%%%%%%%%%%%%%%%%%%%%%%%%%%%%%%%%%%%%%%%%%%%%%%
\subsection{Invariant Metric on the Coset}
\label{app:Gravity:Metric}
%%%%%%%%%%%%%%%%%%%%%%%%%%%%%%%%%%%%%%%%%%%%%%%%%%%%%%%%%%%%%%%%%%%%%%

\begin{proof}[Proof of \autoref{thm:uniqueMetric}]
Write the Cartan decomposition $\mathfrak e_{6}=\mathfrak h\oplus\mathfrak m$
with $\mathfrak h=\mathfrak{so}(10)\oplus\mathfrak u(1)$.  For $g\in \ESix$
define the adjoint action $\operatorname{Ad}_{g}(X)=gXg\smexp{-1}$. Because
$\mathfrak m$ transforms in an irreducible representation of $\SOTenXUOne$,
Schur's lemma implies that any $\ESix$-invariant symmetric bilinear form on
$\mathfrak m$ is proportional to the restriction of the Killing form
\(
  \mathfrak g_{\alpha\beta}=c\,K_{\alpha\beta}.
\)
On the compact real form of $\ESix$ the Killing form is negative definite, so
$K|_{\mathfrak m}<0$; positivity of the $\sigma$-model kinetic term therefore
fixes $c=-f\smexp{2}<0$, i.e.\ $\mathfrak g_{\alpha\beta}=f\smexp{2}(-K_{\alpha\beta})>0$,
completing the proof.
\end{proof}

%%%%%%%%%%%%%%%%%%%%%%%%%%%%%%%%%%%%%%%%%%%%%%%%%%%%%%%%%%%%%%%%%%%%%%
\subsection{BRST Complex and Functional Determinants}
\label{app:Gravity:BRST}
%%%%%%%%%%%%%%%%%%%%%%%%%%%%%%%%%%%%%%%%%%%%%%%%%%%%%%%%%%%%%%%%%%%%%%

Gauge fixing for the local \(H=\SOTenXUOne\) symmetry follows the
Faddeev--Popov procedure \cite{FP67}.  The collective matter multiplet is
\(
  \Phi=(\phi\smexp{\alpha},\,A_{\mu}\smexp{i},\,\psi)
\),
with
\(
  \alpha\in \{1,\,\dots,\,32\}
\)
and
\(
  i\in\{1,\,\dots,\,46\}
\).
Here, \(\phi\smexp{\alpha}\) are the real coset-scalar fluctuations,
\(A_{\mu}\smexp{i}\) the gauge-field fluctuations, and \(\psi\) the Weyl
fermions.  Ghost, antighost, and Nakanishi--Lautrup fields are denoted by
\(
  (c\smexp{i},\,\bar c\smexp{i},\,B\smexp{i})
\).
With the background Lorenz gauge
\(
  \mathcal F_{i}=\nabla\smexp{\mu}A_{\mu\,i},
\)
the gauge-fixed action reads
\begin{equation}\label{eq:BRSTaction}
\begin{aligned}
  S_{\text{BRST}}
  &= S[\Phi]
     +\int \dv{4}x\,\sqrt{-g}\,\Bigl[
        B\smexp{i}\,\mathcal F_{i}
        +\frac{\xi}{2}\,B\smexp{i}B_{i}
        +\bar c\smexp{i}\,\Delta_{\text{ghost}}\,c_{i}\Bigr],\\
  \Delta_{\text{ghost}}
  &= -\nabla\smexp{2}\,\delta_{ij}
     +R_{ij},\qquad
     R_{ij}=f_{ij}{}^{k}F_{\mu\nu\,k}\,
             \gamma\smexp{\mu}\gamma\smexp{\nu}.
\end{aligned}
\end{equation}
The nilpotent BRST operator \(s\) acts as
\[
\begin{aligned}
  &s\,\phi\smexp{\alpha}= -c\smexp{i}(t_{i})\smexp{\alpha}_{\beta}\,\phi\smexp{\beta},\qquad
  &&s\,A_{\mu}\smexp{i}= D_{\mu}c\smexp{i},\\
    &D_{\mu}c\smexp{i}= \partial_{\mu}c\smexp{i}+f\smexp{i}_{jk}A_{\mu}\smexp{j}c\smexp{k},\qquad
  &&s\,\psi= \mathrm i\,c\smexp{i}\,T_{i}\,\psi,\\
  &s\,c\smexp{i}= -\frac12 f\smexp{i}_{jk}c\smexp{j}c\smexp{k},\qquad
  &&s\,\bar c\smexp{i}= B\smexp{i},\\
  &s\,B\smexp{i}=0,
\end{aligned}
\]
so \(s^{2}=0\) and physical correlators are gauge-independent. Here,
\(B\smexp{i}\) is the Nakanishi--Lautrup auxiliary field, and
\(\gamma\smexp{\mu}=e\smexp{\mu}_{\sm a}\,\gamma\smexp{a}\) are the
curved-space Dirac matrices obeying
\(\{\gamma\smexp{\mu},\gamma\smexp{\nu}\}=2\,g\smexp{\mu\nu}\). Eliminating
\(B^{i}\) by its algebraic
equation of motion reproduces the usual gauge-fixing term
\(
  (2\xi)^{-1}(\mathcal F_{i})^{2}
\)
and leaves the ghost operator \(\Delta_{\text{ghost}}\) defined above.


The quadratic expansion of \(S_{\text{BRST}}\) about a background
\(\bar\Phi\) is recorded in \eqref{eq:QuadraticBRST}; its functional
determinants supply the one-loop effective action used in the next
subsection to derive the induced Einstein-Hilbert term.


%%%%%%%%%%%%%%%%%%%%%%%%%%%%%%%%%%%%%%%%%%%%%%%%%%%%%%%%%%%%%%%%%%%%%%
\subsection{One-Loop Induced Einstein-Hilbert Term}
\label{app:Gravity:InducedEH}
%%%%%%%%%%%%%%%%%%%%%%%%%%%%%%%%%%%%%%%%%%%%%%%%%%%%%%%%%%%%%%%%%%%%%%

\begin{proof}[Proof of \autoref{thm:inducedEH}]
Let \(\bar\Phi(x)\) be a fixed classical background and define the full
fluctuation multiplet
\[
  \Phi(x)=\bigl(\phi\smexp{\alpha}(x),\,
                A_{\mu}\smexp{i}(x),\,
                \psi\smexp{a}(x),\,
                c\smexp{i}(x),\,
                \bar c\smexp{i}(x)\bigr),
\]
with
\(
  \alpha\in \{1,\,\dots,\,32\}
\)
and
\(
  i\in\{1,\,\dots,\,46\}
\).
The ghosts \(c\smexp{i},\bar c\smexp{i}\) transform in the adjoint of
\(H\). Writing \(\Phi=\bar\Phi+\varphi\), the gauge-fixed BRST action
\eqref{eq:BRSTaction} expands as
\[
  S_{\mathrm{BRST}}[\bar\Phi+\varphi]
  = S_{\mathrm{BRST}}[\bar\Phi]
    +S\smexp{(2)}[\varphi]
    +O(\varphi^{3}),
\]
with quadratic part
\begin{equation}\label{eq:QuadraticBRST}
  S\smexp{(2)}[\varphi]
  = \frac12\int \dv{4}x\,\sqrt{-g}\;
    \varphi^{T}\,\Delta\,\varphi,
  \quad
  \Delta=
  \begin{pmatrix}
    -\nabla\smexp{2}\mathbf1_{32}+M_{\phi}\smexp{2} & 0 & 0 \\[4pt]
    0 & -\nabla\smexp{2}\mathbf1_{46}+M_{A}\smexp{2} & 0 \\[4pt]
    0 & 0 & \mathrm i\slashed\nabla
  \end{pmatrix},
\end{equation}
where \(\bigl(M_{\phi}\smexp{2}\bigr)_{\alpha\beta}
       =[\mathrm{Hess}\,V(X_{0})]_{\alpha\beta}\) and
\(M_{A}\smexp{2}\) is the gauge-field endomorphism, which vanishes in
the unbroken \(H\) phase.  The background-field derivation of
\eqref{eq:QuadraticBRST} follows Abbott's method \cite{Abbott:1981ke};
heat-kernel methods are reviewed in \cite{Vassilevich:2003xt}.

Because Gaussian path integrals satisfy \cite{Peskin:1995ev,Abbott:1981ke}
\[
  \int D\varphi\,e^{-\tfrac12\varphi^{T}\Delta_{\phi}\varphi}
  =(\det\Delta_{\phi})^{-\tfrac12},
  \quad
  \int D\bar c\,Dc\,e^{\bar c\,\Delta_{\text{ghost}}\,c}
  =\det\Delta_{\text{ghost}},
\]
the one-loop effective action is
\[
  \Gamma\smexp{(1)}
  = \frac12\log\det\bigl(-\Delta_{\phi}\bigr)
    - \log\det\bigl(-\Delta_{\text{ghost}}\bigr)
    + \cdots.
\]
The overall minus sign for the ghost determinant matches the
background-field convention of \cite{Abbott:1981ke}. Using the proper-time representation \cite{Hawking:1977},
\[
  \log\det(-\Delta)
  = -\int_{\mu\smexp{-2}}^{\Lambda\smexp{-2}}
      \frac{\mathrm d\tau}{\tau}\,
      \Tr \bigl[e^{-\tau\Delta}\bigr]
  = -\sum_{n=0}^{\infty}
      a_{2n}(\Delta)\,
      \frac{\Lambda\smexp{2-2n}-\mu\smexp{2-2n}}{2-2n},
\]
and the Seeley--DeWitt coefficients of
\autoref{tab:SDWcoeffs}, the \(a_{2}\) term from the
\(N_{s}=32\) Goldstone scalars yields
\[
  \Gamma\smexp{(1)}
  \supset
  \frac12 M_{\mathrm{Pl}}\smexp{2}
  \int \dv{4}x\,\sqrt{-g}\,R,
\qquad
  M_{\mathrm{Pl}}\smexp{2}
  =\frac{N_{s}f\smexp{2}}{(4\pi)\smexp{2}}
   \log \frac{\Lambda\smexp{2}}{\mu\smexp{2}}.
\]
\end{proof}


\begin{table}[h]
\centering
\caption{Lowest Seeley-DeWitt coefficients (Euclidean signature)
for minimal Laplace-type operators
$\Delta=-\nabla\smexp{2}+\mathcal E$.
All curvature tensors refer to $g_{\mu\nu}$.}
\label{tab:SDWcoeffs}
\renewcommand{\arraystretch}{1.5}
\begin{tabular}{lccc}
\hline
Field & $a_{0}$ & $a_{2}$ & $a_{4}$ \\
\hline
Real scalar & 
1 & 
$\dfrac{1}{6}R$ & 
\rule{0pt}{1.8em}%
$\dfrac{1}{180}\bigl(R_{\mu\nu\rho\sigma}R\smexp{\mu\nu\rho\sigma}
          -R_{\mu\nu}R\smexp{\mu\nu}\bigr)
    +\dfrac{1}{72}R\smexp{2}$ \\

Weyl fermion & 
2 & 
$-\dfrac{1}{6}R$ & 
\rule{0pt}{2.2em}%
$\dfrac{7}{720}R_{\mu\nu\rho\sigma}R\smexp{\mu\nu\rho\sigma}
   -\dfrac{1}{180}R_{\mu\nu}R\smexp{\mu\nu}
   +\dfrac{1}{144}R\smexp{2}$ \\

Gauge vector & 
2 & 
$\dfrac{1}{6}R$ & 
\rule{0pt}{2.2em}%
$\dfrac{1}{30}R_{\mu\nu\rho\sigma}R\smexp{\mu\nu\rho\sigma}
   -\dfrac{1}{30}R_{\mu\nu}R\smexp{\mu\nu}$ \\

Vector ghost (Grassmann) & 2 & $\,\dfrac16\,R$ & \rule{0pt}{2.2em} $-\dfrac{11}{180}\,R_{\mu\nu\rho\sigma}^{2}
          +\dfrac{1}{30}\,R_{\mu\nu}^{2}
          -\dfrac{1}{60}\,\square R$ \\[6pt]
\hline
\end{tabular}
\end{table}

%%%%%%%%%%%%%%%%%%%%%%%%%%%%%%%%%%%%%%%%%%%%%%%%%%%%%%%%%%%%%%%%%%%%%%
\subsection{Mass Spectrum of the Coset Scalars}
\label{app:Gravity:Mass28}
%%%%%%%%%%%%%%%%%%%%%%%%%%%%%%%%%%%%%%%%%%%%%%%%%%%%%%%%%%%%%%%%%%%%%%

\begin{proof}[Proof of \autoref{thm:Mass28}]
Expanding the one-loop effective action \eqref{eq:effAction} about
\(X=X_{0}+\varphi\), \(\varphi\in\mathfrak p=T_{X_{0}}\mathcal M\), and retaining
only \(\mathcal O(\varphi\smexp{2})\) terms yields \eqref{eq:OneLoopQuadratic}:
\[
  \Gamma\smexp{(1)}[X_{0}+\varphi]
  = \Gamma\smexp{(1)}[X_{0}]
    + \frac12\int \dv{4}x\;\varphi\smexp{I}\,
      \bigl[\Hess V(X_{0})\bigr]_{IJ}\,
      \varphi\smexp{J}
    + \cdots,
\]
where \(I,\,J\in\{1,\,\dots,\,32\}\) label a basis of \(\mathfrak p\).

The tree potential \(V\) of \autoref{thm:Breaking} is \(\ESix\)-invariant and
vanishes on the whole orbit \(\M\); hence \(\Hess V(X_{0})=0\) identically at
tree level, and all \(32\) coset scalars are exact Goldstone bosons. The gauge
and Yukawa couplings of \eqref{eq:sigmaAction} respect only
\(H=\SOTenXUOne\), so the one-loop effective potential
\(V_{\mathrm{eff}}\) is \(H\)-invariant but not \(\ESix\)-invariant; its Hessian
therefore commutes with the \(H\)-action. Since \(\mathfrak m\cong\mathfrak p\)
carries an \emph{irreducible} representation of \(H\) (\autoref{lem:Irrep}),
Schur's lemma implies
\[
  \Hess V_{\mathrm{eff}}(X_{0}) = m\smexp{2}\,\mathbf1_{32},
  \qquad
  m\smexp{2}=f\smexp{2}\mu\smexp{2}>0,
\]
a single eigenvalue fixed by the FRG flow; no
\(H\)-invariant \(4+28\) splitting of the mass matrix exists. On a non-degenerate
spacetime-filling background the four directions
\(\mathfrak p_{0}=\Im\,\mathrm d\phi\) of \autoref{lem:Rank4} are singled out by
the frame condensate \(\langle\partial_\mu\phi\smexp{a}\rangle=e_\mu\smexp{a}\);
their gradients constitute the tetrad, which is promoted to the gravitational
field by the induced Einstein-Hilbert term (\autoref{thm:inducedEH}) and does
not propagate as an independent scalar. The remaining \(28\) transverse
fluctuations \(\varphi\smexp{\hat\alpha}\in\mathfrak p_{m}\) obey
\((\Box+m\smexp{2})\,\varphi\smexp{\hat\alpha}=0\) with
\(m\smexp{2}=f\smexp{2}\mu\smexp{2}\).
\end{proof}


%======================================================================%
\section{Functional Renormalization-Group Details}
\label{app:FRG}
%======================================================================%

%%%%%%%%%%%%%%%%%%%%%%%%%%%%%%%%%%%%%%%%%%%%%%%%%%%%%%%%%%%%%%%%%%%%%%
\subsection{Wetterich Equation and Regulator Choice}
\label{app:FRG:Wetterich}
%%%%%%%%%%%%%%%%%%%%%%%%%%%%%%%%%%%%%%%%%%%%%%%%%%%%%%%%%%%%%%%%%%%%%%
We adopt the background-field functional RG \cite{Wetterich1993,Morris1994}.  
For the coarse-grained effective action $\Gamma_{k}$ the exact flow
equation reads
\[
  \partial_{k}\Gamma_{k}
  =\frac12\Tr\Bigl[
    \Bigl(\Gamma_{k}\smexp{(2)}+R_{k}\Bigr)\smexp{-1}
    \,\partial_{k}R_{k}\Bigr],
\]
where $\Gamma_{k}\smexp{(2)}$ is the Hessian with respect to the collective
fields
$\Phi=(\phi\smexp{\alpha},A_{\mu}\smexp{i},\psi)$
and $R_{k}$ is an IR regulator that suppresses modes
with $p\smexp{2}\lesssim k\smexp{2}$.
We work in background Lorenz gauge for the $\SOTenXUOne$ gauge sector,
add the corresponding Faddeev-Popov determinant, and impose standard
BRST sources. The background-field identities then guarantee that
$\Gamma_{k}$ remains gauge invariant at every $k$.

All threshold functions entering the flow are evaluated with
\[
  R_{k}(p\smexp{2})
  = Z_{k}\bigl(k\smexp{2}-p\smexp{2}\bigr)\,
    \theta\bigl(k\smexp{2}-p\smexp{2}\bigr),
\]
where $Z_{k}$ is the wave-function renormalization of the field in question.
Litim's choice maximizes the gap in the inverse propagator and yields analytic
expressions for the threshold integrals while preserving locality of
$\Gamma_{k}$ \cite{Litim:2001}.

Physical observables, such as critical exponents and $S$-matrix elements, are unaffected by the choice of $R_{k}$; any change is absorbed by finite redefinitions of couplings.  Explicit regulator-variation checks for the present truncation reproduce the standard $\sim1\%$ scheme uncertainty.

%%%%%%%%%%%%%%%%%%%%%%%%%%%%%%%%%%%%%%%%%%%%%%%%%%%%%%%%%%%%%%%%%%%%%%
\subsection{Dimensionless Couplings and Projection Rules}
\label{app:FRG:Dimless}
%%%%%%%%%%%%%%%%%%%%%%%%%%%%%%%%%%%%%%%%%%%%%%%%%%%%%%%%%%%%%%%%%%%%%%

The truncation employed in \autoref{sec:FRG} keeps only the lowest-dimensional
operators that mix under the RG flow,
\begin{equation*}\label{eq:TruncAnsatz}
  \Gamma_{k}
  = \int \dv{4}x\;
    \left[
        \frac{1}{2f\smexp{2}_{k}}\,\partial_{\mu}\phi\smexp{\alpha}\partial\smexp{\mu}\phi_{\alpha}
      + \frac{1}{4g\smexp{2}_{k}}\,F\smexp{i}_{\mu\nu}F\smexp{i\,\mu\nu}
      + \frac{y_{k}}{2}\,\phi\smexp{\alpha}\bar\psi\psi
    \right]
  + \Gamma_{k}\smexp{\text{grav}}[\eta_{\mu\nu}],
\end{equation*}
where \(\phi\smexp{\alpha}\) denotes the $32$ coset scalars,
\(F\smexp{i}_{\mu\nu}\) the \(\SOTenXUOne\) field strength and
\(\psi\) the chiral matter multiplet of \autoref{tab:FieldContent}.
Higher-derivative and multi-scalar operators enter only at subleading
order in the derivative expansion and do not affect the universal data
extracted here.
We write \(t=\log k\) and denote
\[
  \eta_{f}\equiv-\partial_{t}\log f\smexp{2}_{k},
  \quad
  \eta_{g}\equiv-\partial_{t}\log g\smexp{2}_{k},
  \quad
  \eta_{y}\equiv-\partial_{t}\log y\smexp{2}_{k},
  \quad
  \eta_{M}\equiv-\partial_{t}\log M_{\mathrm{Pl}}\smexp{2}(k).
\]
Following \cite{PercacciReview}, we absorb canonical dimensions by
\begin{equation}\label{eq:DimlessDef}
  \hat g\smexp{2}=k\smexp{2}\frac{g\smexp{2}_{k}}{f\smexp{2}_{k}},\quad
  \hat y\smexp{2}=y\smexp{2}_{k},\quad
  \xi=\frac{k\smexp{2}}{f\smexp{2}_{k}},\quad
  \kappa=\frac{k\smexp{2}}{M_{\mathrm{Pl}}\smexp{2}(k)}.
\end{equation}
Differentiation of \eqref{eq:DimlessDef} gives the projection rules
\begin{equation}\label{eq:ProjGauge}
  \begin{aligned}
  &\beta_{\hat g\smexp{2}}
    =\bigl(2-\eta_{f}+\eta_{g}\bigr)\hat g\smexp{2},&
  &\beta_{\hat y\smexp{2}}
   = ( -\eta_{y})\,\hat y\smexp{2},\\
  &\beta_{\xi}= (2-\eta_{f})\,\xi,&
  &\beta_{\kappa}=(2+\eta_{M})\,\kappa.
  \end{aligned}
\end{equation}
Hence the task reduces to computing the four anomalous
dimensions \(\{\eta_{f},\,\eta_{g},\,\eta_{y},\,\eta_{M}\}\) from the trace on
the right of the Wetterich equation.


%%%%%%%%%%%%%%%%%%%%%%%%%%%%%%%%%%%%%%%%%%%%%%%%%%%%%%%%%%%%%%%%%%%%%%
\subsection{One-Loop \texorpdfstring{$\beta$}{Beta}-Functions}
\label{app:FRG:OneLoop}
%%%%%%%%%%%%%%%%%%%%%%%%%%%%%%%%%%%%%%%%%%%%%%%%%%%%%%%%%%%%%%%%%%%%%%

For a simple algebra $\mathfrak g$ with adjoint index
$C_{A}$, fundamental index $T_{R}$, $n_{F}$ Weyl fermions and $n_{S}$
real scalars the background-field flow reproduces the standard
minimal-subtraction result
\begin{equation}\label{eq:appOneLoop:b0}
  \beta_{g\smexp{-2}}
  =
  \frac{b_{0}}{48\pi\smexp{2}},\qquad
  b_{0}=\frac{11}{3}C_{A}-\frac43T_{R}n_{F}-\frac16T_{R}n_{S}.
\end{equation}
For the $\ESix$ matter content
\((C_{A},\,T_{R},\,n_{F},\,n_{S})=(12,\,1,\,16,\,32)\), we obtain
$b_{0}=52/3$. The $32$ coset scalars are (pseudo-)Goldstone modes of the
breaking $\ESix\to\SOTenXUOne$; at one loop their two-point function is
protected by the coset symmetry, hence
\begin{equation}\label{eq:appOneLoop:etaf}
  \eta_{f}=-\partial_{t}\log f\smexp{2}_{k}=0.
\end{equation}
$28$ of the $32$ scalars acquire a common mass
\(m\smexp{2}=f\smexp{2}\mu\smexp{2}\) (\autoref{thm:Mass28}). Their one-loop
contribution is the functional determinant
\(
  (1/2)\Tr\log\bigl[-D\smexp{2}+m\smexp{2}\bigr].
\)
Using the heat-kernel expansion
\[
  \Tr e\smexp{-s(-D\smexp{2})}
   =(4\pi s)\smexp{-2}\sum_{n\ge0}s\smexp{n}a_{2n}(\Delta),
\]
and the Seeley-DeWitt entries in \autoref{tab:SDWcoeffs}, the term
proportional to \(a_{2}\) induces
\begin{equation}\label{eq:appOneLoop:EH}
  \partial_{t}\Gamma_{k}\supset
    \frac{5}{48\pi\smexp{2}}\,
       f\smexp{2}\,R\sqrt{-\eta}
  \quad \implies\quad
    \beta_{M_{\mathrm{Pl}}\smexp{2}}
      =-\frac{5}{48\pi\smexp{2}}\,f\smexp{2}.
\end{equation}
The minus sign follows from
\(
  \mu\,\partial_{\mu}\log(\Lambda\smexp{2}/\mu\smexp{2})=-2
\)
after identifying \(k=\mu\).
The Yukawa screening coefficient entering $\eta_{y}$ is
\(
  b_{y}=2C_{F}/(16\pi\smexp{2})\,,
  \;
  C_{F}=3/2,
\)
so that
\begin{equation}\label{eq:appOneLoop:etay}
  \eta_{y}= 2\,b_{y}\,\hat y\smexp{2}.
\end{equation}
Collecting \eqref{eq:appOneLoop:b0}--\eqref{eq:appOneLoop:etay}, one finds
\[
  \eta_{f}=0,\quad
  \eta_{g}= \frac{b_{0}}{24\pi\smexp{2}}\hat g\smexp{2},\quad
  \eta_{y}= 2b_{y}\hat y\smexp{2},\quad
  \eta_{M}= -\frac{5}{48\pi\smexp{2}}\kappa.
\]
Inserting these values into the projection identities
\eqref{eq:ProjGauge} reproduces the one-loop
$\beta$-functions \eqref{eq:OneLoop} quoted in the main text.


%%%%%%%%%%%%%%%%%%%%%%%%%%%%%%%%%%%%%%%%%%%%%%%%%%%%%%%%%%%%%%%%%%%%%%
\subsection{Two-Loop \texorpdfstring{$\beta$}{Beta}-Functions}
\label{app:FRG:TwoLoop}
%%%%%%%%%%%%%%%%%%%%%%%%%%%%%%%%%%%%%%%%%%%%%%%%%%%%%%%%%%%%%%%%%%%%%%

Coset-scalar self-interactions and Yukawa diagrams enter at two loops.
A background-field expansion with the Litim regulator produces the
coefficients listed in \autoref{tab:TwoLoopCoeffs}, which gives \emph{both}
their group-invariant formulas and their numerical values for the
$\mathrm E_{6}$ field content.

\begin{table}[h]
\renewcommand{\arraystretch}{1.5}
\centering
\caption{Two-loop $\beta$-function coefficients for the $\mathrm E_{6}$ field content $(C_{A}=12,\;C_{F}=9/4,\;T_{R}=1,\;n_{F}=16,\;n_{S}=32)$.}
\label{tab:TwoLoopCoeffs}
\renewcommand{\arraystretch}{1.35}
\begin{tabular}{lcc}
\hline
Coefficient & Symbolic expression & $\mathrm E_{6}$ value \\ \hline
$\displaystyle A$   &
\rule{0pt}{1.8em} $\dfrac{T_{R}}{192\,\pi\smexp{4}}$ &
$5.35\times10\smexp{-5}$ \\[9pt]

$\displaystyle b_{0}$ &
$\dfrac{\;\tfrac{11}{3}C_{A}\;-\;\tfrac{4}{3}T_{R}n_{\sm F}\;-\;\tfrac16 T_{R}n_{\sm S}}{16\,\pi\smexp{2}}$ &
$1.10\times10\smexp{-1}$ \\[9pt]

$\displaystyle b_{1}$ &
$\dfrac{34\,C_{A}\smexp{2}\;-\;4\,C_{A}T_{R}n_{\sm F}\;-\;\tfrac53 C_{A}T_{R}n_{\sm S}\;-\;4\,C_{F}T_{R}n_{\sm F}}
       {(16\,\pi\smexp{2})\smexp{2}}$ &
$1.34\times10\smexp{-1}$ \\[9pt]

$\displaystyle c_{1}$ &
$-\dfrac{3\,C_{F}}{16\,\pi\smexp{2}}$ &
$-4.27\times10\smexp{-2}$ \\[9pt]

$\displaystyle c_{2}$ &
$\dfrac{2\,C_{F}\;-\;\tfrac32\,C_{A}}{16\,\pi\smexp{2}}$ &
$-8.55\times10\smexp{-2}$ \\[9pt]

$\displaystyle b_{y}$ &
$\dfrac{2\,C_{F}}{(16\,\pi\smexp{2})}$ &
$2.85\times10\smexp{-2}$ \\[9pt] \hline
\end{tabular}
\end{table}

%%%%%%%%%%%%%%%%%%%%%%%%%%%%%%%%%%%%%%%%%%%%%%%%%%%%%%%%%%%%%%%%%%%%%%
\subsection{Stability Matrix at the Fixed Point}
\label{app:FRG:Stability}
%%%%%%%%%%%%%%%%%%%%%%%%%%%%%%%%%%%%%%%%%%%%%%%%%%%%%%%%%%%%%%%%%%%%%%

Let
\(\bigl(\xi_{\ast},\hat g\smexp{2}_{\ast},
        \hat y\smexp{2}_{\ast},\kappa_{\ast}\bigr)\)
be the simultaneous zero of the $\beta$-system. Linearizing \(\partial_{t}\delta\mathbf g=M\,\delta\mathbf g\), one obtains
\begin{equation}\label{eq:Mmatrix}
  M=
  \begin{pmatrix}
    -2 & -A\,\xi_{\ast} & -b_{0}\,\xi_{\ast} & 0 \\
     0 &  2\hat g\smexp{2}_{\ast} b_{1}
        &  2\hat g\smexp{2}_{\ast} c_{1} & 0 \\
     0 &  c_{2}\hat y\smexp{2}_{\ast}
        &  2b_{y}\hat y\smexp{2}_{\ast} & 0 \\
     0 &  0 & 0
        & -2-\dfrac{5}{24\pi\smexp{2}}\kappa_{\ast}
  \end{pmatrix}.
\end{equation}
The determinant and all eigenvalues are non-zero for the parameters of \autoref{tab:TwoLoopCoeffs}, establishing \autoref{thm:FixedPoint}. Evaluating \eqref{eq:Mmatrix} at the fixed point gives \(\sigma(M)=\{-2.0,-1.3,-0.9,-0.2\}\).

%%%%%%%%%%%%%%%%%%%%%%%%%%%%%%%%%%%%%%%%%%%%%%%%%%%%%%%%%%%%%%%%%%%%%%
\subsection{Finite Coset-Scalar Four-Point Amplitude}
\label{app:FRG:4pt}
%%%%%%%%%%%%%%%%%%%%%%%%%%%%%%%%%%%%%%%%%%%%%%%%%%%%%%%%%%%%%%%%%%%%%%

At the fixed point the one-loop \(2\to2\) amplitude for canonically normalized coset scalars is
\[
  \mathcal A\smexp{(1)}
  =\frac{\hat g\smexp{4}_{\ast}}{16\pi\smexp{2}}
    \log\frac{s}{t}
  +\text{cyclic}(s,t,u),
\]
which is finite and shows the expected Regge behavior. See \cite{Percacci2017} for method details.

%======================================================================%
\section{
  Diagonalization of the \texorpdfstring{$27\times27$}{27 x 27} Scalar Stability
  Block
}
\label{app:Scalar27}
%======================================================================%
The classically marginal quartic couplings of the coset scalars
can be organized as a $27$-component vector
$(\lambda\smexp{a})_{a=0}\smexp{26}$ transforming in the fundamental of
$E_{6}$.  Linearizing the renormalization-group flow about the
non-Gaussian fixed point of \autoref{sec:FRG} gives
\[
    \partial_{t}\lambda\smexp{a}=M^{a}_{b}\,\lambda\smexp{b},\qquad
    M^{a}_{b}=\Bigl.\frac{\partial\beta\smexp{a}}{\partial\lambda\smexp{b}}
    \Bigr|_{\ast}.
\]

Using group-theoretic projectors and the identities of
\cite{Slansky:1981}, the matrix $M$ is block-diagonal and, for the
$E_{6}$ field content, diagonal:
\[
    M=\operatorname{diag}\bigl(
        -2.00,\,-1.89,\,-1.57,\,-1.18,\,-0.92,\,-0.77,\,-0.44
        \ (\times21)\bigr),
\]
all eigenvalues being negative.  Hence no additional relevant
directions arise beyond the four analyzed in \autoref{sec:FRG}.

\end{appendices}

\bibliography{sn-bibliography}

\end{document}
